% Optimization Methods and Software submission
% Taylor & Francis Interact layout (reference style S)
% Ported from the Elsevier cas-sc manuscript (optimization_ready.tex).
\documentclass[]{interact}
\usepackage{fontspec}
\setmainfont{Latin Modern Roman}[SmallCapsFont={Latin Modern Roman Caps}]

\usepackage{epstopdf}% To incorporate .eps illustrations using PDFLaTeX, etc.
\usepackage[caption=false]{subfig}% Support for small, `sub' figures and tables

\usepackage[numbers,sort&compress]{natbib}% Citation support using natbib.sty
\bibpunct[, ]{[}{]}{,}{n}{,}{,}% Citation support using natbib.sty
\renewcommand\bibfont{\fontsize{10}{12}\selectfont}% Bibliography support using natbib.sty

% --- Packages carried over from the source manuscript ---
\usepackage{amssymb}
\usepackage{float}
\usepackage{algorithm}
\usepackage{multirow}
\usepackage{array}
\usepackage{longtable}
\usepackage{enumitem}
\usepackage{xcolor,url}\graphicspath{{figures/}}
\restylefloat{figure}

% --- Theorem-like environments ---
\theoremstyle{plain}
\newtheorem{theorem}{Định lý}[section]
\newtheorem{lemma}[theorem]{Bổ đề}
\newtheorem{corollary}[theorem]{Hệ quả}
\newtheorem{proposition}[theorem]{Mệnh đề}

\theoremstyle{definition}
\newtheorem{definition}[theorem]{Định nghĩa}
\newtheorem{example}[theorem]{Ví dụ}
\newtheorem{assumption}{Giả thiết}

\theoremstyle{remark}
\newtheorem{remark}{Nhận xét}
\newtheorem{notation}{Ký hiệu}

% --- Author macros carried over from the source manuscript ---
\newcommand{\algstep}[1]{\smallskip\noindent\textbf{#1}}
\newcommand{\R}{\mathbb{R}}
\newcommand{\N}{\mathbb{N}}
\newcommand{\norm}[1]{\left\lVert #1 \right\rVert}
\newcommand{\abs}[1]{\left\lvert #1 \right\rvert}
\newcommand{\inner}[2]{\left\langle #1,\, #2 \right\rangle}
\DeclareMathOperator*{\argmin}{arg\,min}
\DeclareMathOperator*{\argmax}{arg\,max}
\DeclareMathOperator{\conv}{conv}
\DeclareMathOperator{\diag}{diag}
\DeclareMathOperator{\vecop}{vec}
\newcommand{\cvx}{\mathrm{cvx}}
\newcommand{\ps}{\mathrm{ps}}
\newcommand{\dtap}{\textsc{DTAP}}

% --- Vietnamese names for the class's fixed labels ---
\renewcommand\abstractname{TÓM TẮT}
\renewcommand\keywordsname{TỪ KHÓA}
\renewcommand\amscodename{PHÂN LOẠI AMS}
\renewcommand\tablename{Bảng}
\renewcommand\figurename{Hình}
\renewcommand\refname{Tài liệu tham khảo}
\renewcommand\proofname{Chứng minh}
\floatname{algorithm}{Thuật toán}
\makeatletter
\gdef\@history{\bfseries{BẢN DỊCH\\}\par}
\gdef\@received{Bản dịch tiếng Việt của \texttt{DTAP Finally v2/main\_document.tex}, chỉ dùng để rà soát nội dung, không nộp. Công thức, nhãn và cách đánh số giữ nguyên như bản tiếng Anh.}
\makeatother

\begin{document}

\title{Thuật toán Chiếu Thích nghi Hai Nhánh cho một bài toán chấp nhận tách phi tuyến với các ràng buộc hỗn hợp lồi--giả lồi}

\author{
\name{Nguyen Van The\textsuperscript{a}\thanks{Email: the.nv227153@sis.hust.edu.vn}, Tran Ngoc Thang\textsuperscript{a}\thanks{Email: thang.tranngoc@hust.edu.vn} và Nguyen Hong Son\textsuperscript{a}\thanks{LIÊN HỆ Nguyen Hong Son. Email: Son.NH240066D@sis.hust.edu.vn}}
\affil{\textsuperscript{a}Khoa Toán -- Tin, Đại học Bách khoa Hà Nội, số~1 Đại Cồ Việt, phường Bạch Mai, Hà Nội 100000, Việt Nam}
}

\maketitle

\begin{abstract}
Chúng tôi nghiên cứu một bài toán chấp nhận tách phi tuyến có tập ảnh là một hộp chuẩn, trong đó các hàm thành phần khả vi và mỗi hàm là lồi hoặc giả lồi. Chúng tôi đề xuất thuật toán Chiếu Thích nghi Hai Nhánh (DTAP): thuật toán chia các ràng buộc thành một nhánh lồi và một nhánh giả lồi, và trên mỗi nhánh làm việc với các ràng buộc gần như bị vi phạm nhiều nhất. DTAP kết hợp một bước kiểu Polyak điều chỉnh theo phần dư trên nhánh lồi với một bước giảm dần có bảo vệ trên nhánh giả lồi, thông qua một trọng số tỉ lệ thay đổi theo phần dư, rồi chiếu lên tập lồi cơ sở. Các cỡ bước của thuật toán không đòi hỏi biết trước chuẩn toán tử, hằng số Lipschitz hay mô-đun H\"older; chứng minh sử dụng các cận địa phương trên miền bị chặn chứa các điểm lặp. Với giả thiết nhất quán và khả vi liên tục, chúng tôi chứng minh rằng toàn bộ dãy lặp hội tụ mạnh tới một điểm chấp nhận được. Làm ứng dụng, chúng tôi đưa bài toán chọn đặc trưng về dạng bài toán trên, với các ràng buộc được xây dựng từ các tiêu chí chọn đặc trưng dạng phân thức và dạng tập thô mờ. Thực nghiệm trên các bộ dữ liệu UCI cho thấy các đặc trưng được chọn giữ độ chính xác phân loại gần với độ chính xác khi dùng toàn bộ đặc trưng, còn các bài toán thử có ràng buộc giả lồi không lồi cho phép kiểm tra nhánh giả lồi; chúng tôi cũng so sánh DTAP với các phương pháp dưới gradient và khảo sát độ nhạy của DTAP đối với điểm xuất phát và thang cỡ bước.
\end{abstract}

\begin{keywords}
Bài toán chấp nhận tách; Tối ưu giả lồi; Vô hướng hóa Tchebyshev; Cỡ bước thích nghi; Hội tụ mạnh; Chọn đặc trưng
\end{keywords}
\begin{amscode}
47J25; 90C25; 90C26; 65K05; 68T05
\end{amscode}

\medskip
\noindent\fbox{\parbox{0.97\textwidth}{\footnotesize
\textbf{Bảng thuật ngữ Anh--Việt dùng trong bản dịch.}
split feasibility problem (SFP): bài toán chấp nhận tách;\
feasible set / point: tập / điểm chấp nhận được;\
track: nhánh;\
normal box: hộp chuẩn;\
step size: cỡ bước;\
(approximate) active set: tập active (xấp xỉ);\
violation: mức vi phạm;\
residual: phần dư;\
residual-scaled step: bước điều chỉnh theo phần dư;\
safeguarded diminishing step: bước giảm dần có bảo vệ;\
proportional weight: trọng số tỉ lệ;\
consistency: tính nhất quán;\
pseudoconvex / quasiconvex / quasilinear: giả lồi / tựa lồi / tựa tuyến tính;\
sublevel set: tập mức dưới;\
cluster point: điểm tụ;\
metric projection: phép chiếu mêtric;\
quasi-Fej\'er: tựa Fej\'er;\
subgradient method: phương pháp dưới gradient;\
tuned: đã tinh chỉnh;\
tolerance: sai số cho phép;\
budget: số vòng lặp (hoặc số điểm) tối đa;\
evaluated point: điểm đã tính;\
slack: độ chênh $b_i-F_i(w)$;\
witness: điểm minh chứng;\
fold, seed, mask: fold, seed, mặt nạ (giữ nguyên tên tiếng Anh).}}
\medskip

\section{Giới thiệu}\label{sec:intro}

Bài toán chấp nhận tách (SFP) được Censor và Elfving~\citep{censor1994} đưa ra như một mô hình cho các bài toán ngược có ràng buộc. Cho các tập lồi đóng khác rỗng $C\subseteq\R^n$, $Q\subseteq\R^m$ và một ma trận $\mathcal A\in\R^{m\times n}$, bài toán yêu cầu
\begin{equation}\label{eq:sfp-intro}
\text{tìm }x^\star\in C\quad\text{sao cho}\quad\mathcal A x^\star\in Q.
\end{equation}
Thuật toán $CQ$ của Byrne~\citep{byrne2002} lặp theo công thức
\begin{equation}\label{eq:cq-intro}
x^{k+1}=P_C\bigl(x^k-\tau\mathcal A^\top(I-P_Q)\mathcal A x^k\bigr),
\qquad 0<\tau<2/\norm{\mathcal A}^2.
\end{equation}
Ở đây $\norm{\mathcal A}$ là chuẩn phổ của $\mathcal A$. Các biến thể nới lỏng chiếu lên các nửa không gian tựa~\citep{yang2004}; các cỡ bước tự thích nghi tránh việc phải ước lượng $\norm{\mathcal A}$~\citep{zhang2009adaptive,zhao2011adaptive,lopez2012}; phương pháp tách toán tử~\citep{zhang2012simultaneous} và phương pháp hướng luân phiên~\citep{yang2014alternating} phân rã bài toán. Các ý tưởng này đã quen thuộc; DTAP vận dụng chúng cho trường hợp phi tuyến.

Thay $\mathcal A$ bởi một ánh xạ phi tuyến $F$, ta được SFP phi tuyến~\citep{li2013nsfp,yen2019}
\begin{equation}\label{eq:nsfp-intro}
\text{tìm }x^\star\in K\quad\text{sao cho}\quad F(x^\star)\in Q,
\end{equation}
trong đó $K\subseteq\R^n$ là tập khác rỗng, đóng và lồi, $Q=\prod_{i=1}^m(-\infty,b_i]$ là một hộp chuẩn, và mỗi $F_i$ khả vi, lồi hoặc giả lồi trên $K$. Vì các hàm giả lồi có tập mức dưới lồi, tập chấp nhận được $\Omega=\{x\in K:F(x)\in Q\}$ là tập lồi. Ta tìm một điểm thuộc $\Omega$, chứ không tìm điểm cực tiểu của một hàm $F_i$ nào.

Vì $Q$ là một hộp, khoảng cách $\ell_\infty$ từ $F(x)$ tới $Q$ có dạng hiển
\[
J(x)=d_\infty(F(x),Q)=\max_{1\le i\le m}(F_i(x)-b_i)^+.
\]

Li và cộng sự~\citep{li2013nsfp} giả thiết hàm lân cận tổng hợp là lồi, điều mà tính giả lồi của từng thành phần không bảo đảm. Yen và cộng sự~\citep{yen2019} dùng cùng độ đo vi phạm kiểu max và cùng kỹ thuật tập active, nhưng đòi hỏi các thành phần tựa tuyến tính, một lớp hàm không so sánh được với lớp hàm của bài này. Hu và cộng sự~\citep{hu2022quasi} xét các bất đẳng thức tựa lồi, một lớp chứa lớp hàm của bài này, nhưng không sử dụng bất đẳng thức gradient của các thành phần lồi.

Ràng buộc lồi và ràng buộc giả lồi cần những bước đi khác nhau. Với một ràng buộc lồi, bất đẳng thức gradient chặn dưới $\inner{\nabla F_i(x)}{x-x^\star}$, với mọi $x^\star\in\Omega$, bởi mức vi phạm, điều này biện minh cho một bước kiểu Polyak điều chỉnh theo phần dư~\citep{polyak1969,lopez2012}; với một ràng buộc giả lồi, ta chỉ biết dấu của tích vô hướng này, nên dùng một bước giảm dần có bảo vệ. DTAP thực hiện hai bước này trên hai nhánh riêng và kết hợp chúng bằng một trọng số tỉ lệ; chứng minh chỉ dùng trọng số này để bảo đảm một nhánh bị vi phạm kéo dài vẫn còn trọng số dương (Nhận xét~\ref{rem:weight-scope}).

Các đóng góp của bài như sau.
\begin{enumerate}[label=\textbf{C\arabic*.},leftmargin=2.4em,itemsep=3pt]
\item Một phương pháp hai nhánh dùng bất đẳng thức gradient trên các ràng buộc lồi và chỉ dùng thông tin về dấu trên các ràng buộc giả lồi, không cần giả thiết hàm lân cận lồi như~\citep{li2013nsfp}, cũng không cần mô-đun H\"older như cỡ bước động của~\citep{hu2022quasi}.
\item Sự hội tụ của toàn bộ dãy lặp tới một điểm chấp nhận được dưới Giả thiết~\ref{ass:standing} (Định lý~\ref{thm:strong}).
\item Một ứng dụng chọn đặc trưng xây dựng từ các tiêu chí dạng phân thức và dạng tập thô mờ, có các ràng buộc lồi (Mệnh đề~\ref{prop:fs-reg}), cùng các bài toán thử có ràng buộc giả lồi không lồi; thực nghiệm báo cáo mức hoạt động của các ràng buộc, so sánh DTAP với các phương pháp dưới gradient đã tinh chỉnh, và đo độ nhạy của DTAP đối với điểm xuất phát và thang cỡ bước.
\end{enumerate}

Mục~\ref{sec:prelim} trình bày kiến thức chuẩn bị. Mục~\ref{sec:convergence} trình bày \dtap{} và phân tích hội tụ. Mục~\ref{sec:experiments} báo cáo các thực nghiệm, và Mục~\ref{sec:conclusion} kết luận.


\section{Kiến thức chuẩn bị}\label{sec:prelim}
Trong toàn bài, $H=\R^n$ là không gian Euclid với tích vô hướng $\inner{\cdot}{\cdot}$ và chuẩn sinh bởi tích vô hướng $\norm{\cdot}$. Mọi phát biểu về hội tụ đều hiểu theo tô-pô chuẩn. Mọi dãy bị chặn đều có dãy con hội tụ. Các kết quả quen thuộc về phép chiếu và về lồi suy rộng dưới đây có thể tìm trong~\citep{bauschke2011,mangasarian1994}.

Với một tập lồi đóng khác rỗng $K\subseteq H$, phép chiếu mêtric $P_K(x)$ là điểm cực tiểu duy nhất của $\norm{y-x}$ trên $y\in K$.
\begin{proposition}\label{prop:proj-char}
Với $x\in H$ và $p\in K$, ta có $p=P_K(x)$ khi và chỉ khi
\begin{equation}\label{eq:proj-char}
\inner{x-p}{y-p}\le0\qquad\text{với mọi }y\in K.
\end{equation}
\end{proposition}
Đặc trưng này dẫn tới Bổ đề~\ref{lem:proj-nonexp} ở dưới.

Hàm $f:K\to\R$ được gọi là lồi nếu
$f((1-t)x+ty)\le(1-t)f(x)+tf(y)$ với $x,y\in K$ và $t\in[0,1]$.
Điều kiện đạo hàm yếu hơn sau đây do Mangasarian~\citep{mangasarian1965} đưa ra.
\begin{definition}\label{def:pseudoconvex}
Hàm $f$ khả vi trên một lân cận mở của $K$ được gọi là giả lồi trên $K$ nếu với mọi $x,y\in K$,
\begin{equation}\label{eq:contrapositive}
f(y)<f(x)\quad\Longrightarrow\quad\inner{\nabla f(x)}{y-x}<0.
\end{equation}
\end{definition}
Mọi hàm lồi khả vi đều giả lồi. Hàm giả lồi có các tập mức dưới lồi, mặc dù bản thân hàm có thể không lồi.

\section{Thuật toán DTAP và phân tích hội tụ}\label{sec:convergence}
Ta phát biểu NSFP với hộp chuẩn, trình bày DTAP, và chứng minh toàn bộ dãy lặp hội tụ
tới một điểm chấp nhận được.

\subsection{Phát biểu bài toán và vô hướng hóa Tchebyshev}\label{ssec:model-stmt}
Cho $K\subseteq H$ là tập lồi đóng khác rỗng, và $U\subseteq H$ là một lân cận mở
của $K$. Cố định số nguyên $m\ge1$, cho $F=(F_1,\dots,F_m):U\to\R^m$ và
$b=(b_1,\dots,b_m)\in\R^m$, và đặt hộp chuẩn
$Q:=(-\infty,b_1]\times\cdots\times(-\infty,b_m]$. Ta xét trường hợp sau
của~\eqref{eq:nsfp-intro}:
\begin{equation}\label{eq:nsfp-box}
\text{tìm } x^\star\in\Omega,\qquad
\Omega:=\bigl\{x\in K: F_i(x)\le b_i,\ i=1,\dots,m\bigr\},
\end{equation}
và gọi $\Omega$ là \emph{tập chấp nhận được}. Để biến điều kiện $F(x)\in Q$ thành một hàm đánh giá
dễ xử lý, ta đo phần dư theo chuẩn sup. Với $z=F(x)$,
$d_\infty(z,Q)=\inf_{q\le b}\max_{1\le i\le m}\abs{z_i-q_i}$. Do cấu trúc hộp,
cận dưới đúng tách được theo từng tọa độ: nếu $z_i\le b_i$ thì
$\inf_{q_i\le b_i}\abs{z_i-q_i}=0$ (đạt tại $q_i=z_i$), còn nếu $z_i>b_i$ thì nó bằng
$z_i-b_i$ (đạt tại $q_i=b_i$). Do đó $\inf_{q_i\le b_i}\abs{z_i-q_i}=(z_i-b_i)^+$ với
$(t)^+:=\max\{t,0\}$, nên
\begin{equation}\label{eq:dinfty}
d_\infty(z,Q)=\max_{1\le i\le m}(z_i-b_i)^+ .
\end{equation}
Định nghĩa \emph{hàm sai số}
\begin{equation}\label{eq:Jdef}
J(x):=d_\infty\bigl(F(x),Q\bigr)=\max_{1\le i\le m}\bigl(F_i(x)-b_i\bigr)^+ .
\end{equation}
Khi đó $x\in K$ là điểm chấp nhận được khi và chỉ khi $J(x)=0$; khi $\Omega\ne\emptyset$, các điểm
chấp nhận được chính là các điểm cực tiểu của $J$ trên $K$. Công thức~\eqref{eq:Jdef} cho một phần dư dạng hiển của tính chấp nhận được, và các điểm
đạt max của nó chỉ ra những ràng buộc bị vi phạm nhiều nhất. Mục~\ref{ssec:dual-track}
đưa vào một tập active xấp xỉ để xử lý tính không trơn của hàm max.

\subsection{Phân rã thành hai nhánh}\label{ssec:dual-track}
Ta chia $I:=\{1,\dots,m\}$ thành \emph{nhánh lồi} $I^{\cvx}$ và
\emph{nhánh giả lồi} $I^{\ps}$, mỗi nhánh có quy tắc bước riêng. Nhánh lồi dùng bước Polyak~\citep{polyak1969};
nhánh giả lồi dùng bước giảm dần có bảo vệ.
Một trọng số tỉ lệ kết hợp hai bước cập nhật này.

Cố định một tham số $\delta>0$. Với $x\in U$, định nghĩa mức vi phạm của từng thành phần
\begin{equation}\label{eq:vi}
v_i(x):=\bigl(F_i(x)-b_i\bigr)^+,
\end{equation}
khi đó $J(x)=\max_i v_i(x)$ theo~\eqref{eq:Jdef}. Giả thiết~\ref{ass:standing} phát biểu chính xác
phép chia này. \emph{Mức vi phạm lớn nhất trên từng nhánh} là
\begin{align}
J^{\cvx}(x)&:=\max_{i\in I^{\cvx}}v_i(x)\quad(\,{=}\,0\text{ nếu }I^{\cvx}=\emptyset),\label{eq:Jcvx}\\
J^{\ps}(x)&:=\max_{j\in I^{\ps}}v_j(x)\quad(\,{=}\,0\text{ nếu }I^{\ps}=\emptyset),\label{eq:Jps}
\end{align}
nên $J(x)=\max\{J^{\cvx}(x),J^{\ps}(x)\}$. Để xử lý tính không trơn của hàm max,
trên mỗi nhánh ta thay tập các chỉ số đạt giá trị lớn nhất bằng một \emph{tập active xấp xỉ}
được nới rộng một lượng $\delta$ (``active'' nghĩa là mức vi phạm cách giá trị lớn nhất
không quá $\delta$):
\begin{align}
K^{\cvx}(x)&:=\begin{cases}\{i\in I^{\cvx}:v_i(x)\ge J^{\cvx}(x)-\delta\} & \text{nếu }J^{\cvx}(x)>0,\\ \emptyset & \text{nếu }J^{\cvx}(x)=0,\end{cases}\label{eq:Kcvx}\\[2pt]
K^{\ps}(x)&:=\begin{cases}\{j\in I^{\ps}:v_j(x)\ge J^{\ps}(x)-\delta\} & \text{nếu }J^{\ps}(x)>0,\\ \emptyset & \text{nếu }J^{\ps}(x)=0.\end{cases}\label{eq:Kps}
\end{align}
Trên các tập này ta lập \emph{các hàm và các hướng địa phương theo từng nhánh}
\begin{align}
\Phi^{\cvx}(x)&:=\tfrac12\!\!\sum_{i\in K^{\cvx}(x)}\!\! v_i(x)^2, & g^{\cvx}(x)&:=\!\!\sum_{i\in K^{\cvx}(x)}\!\! v_i(x)\nabla F_i(x),\label{eq:phi-g-cvx}\\
& & g^{\ps}(x)&:=\!\!\sum_{j\in K^{\ps}(x)}\!\! v_j(x)\nabla F_j(x).\label{eq:phi-g-ps}
\end{align}
Các vector $g^{\cvx}(x)$, $g^{\ps}(x)$ là hai hướng tìm kiếm, còn $\Phi^{\cvx}(x)$
là tử số trong bước Polyak trên nhánh lồi. Dọc theo dãy lặp, ta viết tắt
$J^{\cvx}_k:=J^{\cvx}(x^k)$, $g^{\cvx}_k:=g^{\cvx}(x^k)$, $\Phi^{\cvx}_k:=\Phi^{\cvx}(x^k)$,
và tương tự trên nhánh giả lồi.

\subsection{Các giả thiết}\label{ssec:assumptions}
\begin{assumption}\label{ass:standing}
Ta giả thiết \textup{(A1)--(A3)}:
\begin{itemize}[leftmargin=2.4em,itemsep=2pt]
\item[\textup{(A1)}] \textup{(a)} $\Omega\ne\emptyset$ (tính nhất quán). \textup{(b)} $I=\{1,\dots,m\}$
được chia thành $I=I^{\cvx}\cup I^{\ps}$; mỗi $F_i:U\to\R$ lồi trên $K$ nếu
$i\in I^{\cvx}$ và giả lồi trên $K$ nếu $i\in I^{\ps}$.
\item[\textup{(A2)}] Mỗi $F_i:U\to\R$ khả vi liên tục trên $U$.
\item[\textup{(A3)}] \textup{(i)} $\delta>0$ cố định. \textup{(ii)}
$\rho_{\cvx}\in(0,2)$. \textup{(iii)} $\{\gamma_k\}\subset(0,\infty)$ thỏa
$\sum_{k=0}^\infty\gamma_k^2<\infty$ và $\sum_{k=0}^\infty\gamma_k=\infty$.
\end{itemize}
\end{assumption}
Việc phát biểu trên một lân cận mở cho phép các hàm phân thức có mẫu số
dương trên $K$; thuật toán chỉ tính giá trị hàm và gradient tại các điểm thuộc $K$.
Dù một số hàm thành phần có thể không lồi, tập $\Omega$ vẫn lồi: các hàm giả lồi khả vi
là tựa lồi trên $K$, nên mọi tập mức dưới của các ràng buộc trong $K$ đều lồi.
\begin{example}\label{ex:gamma}
Dãy $\gamma_k:=1/(k+1)$ thỏa (A3)(iii), vì $\sum_{k\ge0}(k+1)^{-2}<\infty$
trong khi $\sum_{k\ge0}(k+1)^{-1}=\infty$. Tổng quát hơn, $\gamma_k=(k+c)^{-s}$ với $c>0$
và $s\in(1/2,1]$ đều thỏa mãn; các thực nghiệm dùng $s=0.55$.
\end{example}

\subsection{Thuật toán DTAP}\label{ssec:algorithm}
Thuật toán Chiếu Thích nghi Hai Nhánh kết hợp hai cỡ bước riêng của hai nhánh
thông qua trọng số tỉ lệ $\mu_k$, rồi chiếu bước cập nhật tổng hợp lên tập lồi cơ sở
$K$. Thuật toán được phát biểu trong Thuật toán~\ref{alg:dtap}.

\begin{algorithm}[H]
\caption{Chiếu Thích nghi Hai Nhánh (DTAP) cho NSFP với hộp chuẩn
\eqref{eq:nsfp-box}.}
\label{alg:dtap}
\small
\noindent\textbf{Đầu vào:} ánh xạ $F=(F_1,\dots,F_m)$, các ngưỡng $b$, tập cơ sở $K$, phép chia
chỉ số $I^{\cvx}\cup I^{\ps}$.\\
\textbf{Tham số:} $\delta>0$, sai số cho phép $\varepsilon>0$, $\rho_{\cvx}\in(0,2)$,
dãy $\{\gamma_k\}$ thỏa (A3)(iii).\\
\textbf{Khởi tạo:} chọn $x^0\in K$; đặt $k\leftarrow 0$.

\algstep{Bước 1 (các đại lượng theo từng nhánh).} Nếu $J(x^k)\le\varepsilon$ thì dừng và trả về $x^k$; nếu không:
\begin{itemize}[leftmargin=1.6em,itemsep=1pt,topsep=2pt]
\item[(1a)] Tính $v_i(x^k)$ với mọi $i$, và $J^{\cvx}_k=\max_{i\in I^{\cvx}}v_i(x^k)$,
$J^{\ps}_k=\max_{j\in I^{\ps}}v_j(x^k)$ (bằng $0$ nếu nhánh tương ứng rỗng).
\item[(1b)] Lập các tập active xấp xỉ $K^{\cvx}_k,K^{\ps}_k$ theo
\eqref{eq:Kcvx}--\eqref{eq:Kps}.
\item[(1c)] Lập $\Phi^{\cvx}_k$, $g^{\cvx}_k$, $g^{\ps}_k$ theo
\eqref{eq:phi-g-cvx}--\eqref{eq:phi-g-ps}.
\item[(1d)] Cỡ bước trên từng nhánh (bước Polyak trên nhánh lồi, bước có bảo vệ trên
nhánh giả lồi):
\[
\eta^{\cvx}_k=\begin{cases}\rho_{\cvx}\,\dfrac{2\Phi^{\cvx}_k}{\norm{g^{\cvx}_k}^2} & \text{nếu }\norm{g^{\cvx}_k}>0,\\ 0 & \text{ngược lại,}\end{cases}
\qquad
\eta^{\ps}_k=\frac{\gamma_k}{\max\{1,\norm{g^{\ps}_k}\}} .
\]
\item[(1e)] Trọng số tỉ lệ:
$\displaystyle \mu_k=\frac{J^{\cvx}_k}{J^{\cvx}_k+J^{\ps}_k}$ nếu
$J^{\cvx}_k+J^{\ps}_k>0$, và $\mu_k=1/2$ trong trường hợp còn lại.
\item[(1f)] Kết hợp hai hướng:
$\displaystyle y^k=x^k-\bigl[\mu_k\eta^{\cvx}_k g^{\cvx}_k+(1-\mu_k)\eta^{\ps}_k g^{\ps}_k\bigr].$
\end{itemize}
\algstep{Bước 2 (chiếu).}\quad $x^{k+1}=P_K(y^k)$; đặt $k\leftarrow k+1$ và quay lại Bước~1.
\end{algorithm}

\begin{remark}\label{rem:weight-scope}
Trọng số tỉ lệ là một cách phân bổ thích nghi dựa trên hai mức vi phạm lớn nhất;
chúng tôi không khẳng định đó là lựa chọn duy nhất làm DTAP hội tụ, cũng không khẳng định nó tối ưu.
Dạng cụ thể của trọng số được dùng trong~\eqref{eq:p10} để kiểm soát phần dư trên nhánh lồi, và trong
Khẳng định~2 thuộc Giai đoạn~2 của Định lý~\ref{thm:strong} để bảo đảm một nhánh giả lồi
bị vi phạm kéo dài vẫn giữ trọng số dương. Mặc dù bất đẳng thức một bước trong Bổ đề~\ref{lem:fejer} đúng
với mọi $\mu_k\in[0,1]$, riêng bất đẳng thức này không suy ra sự hội tụ với mọi dãy trọng số.
\end{remark}

\begin{remark}\label{rem:regularity-scope}
Phép lặp không đòi hỏi biết trước chuẩn toán tử, hằng số Lipschitz hay
mô-đun H\"older. Tuy vậy, chứng minh có dùng các cận hữu hạn của giá trị hàm và
gradient trên một tập compact chứa các điểm lặp, sau khi đã chứng minh tính bị chặn.
Gradient bị chặn kéo theo các hàm là Lipschitz trên tập compact lồi đó,
nhưng các cận này không phải là đầu vào của thuật toán. Giả thiết (A2) không
kéo theo tính Lipschitz của gradient, và định lý cũng không cần tính chất đó.
Người dùng vẫn phải chọn $\rho_{\cvx}$, $\delta$ và $\{\gamma_k\}$; các tham số này không
phụ thuộc vào hằng số nào của bài toán.
\end{remark}

Để phân tích, ta viết số gia tổng hợp dưới dạng
\begin{equation}\label{eq:dk}
d_k:=\mu_k\eta^{\cvx}_k g^{\cvx}_k+(1-\mu_k)\eta^{\ps}_k g^{\ps}_k,
\end{equation}
nên $y^k=x^k-d_k$ và $x^{k+1}=P_K(x^k-d_k)$.
\subsection{Chi phí mỗi vòng lặp}\label{ssec:complexity}
Mỗi vòng lặp của \dtap{} tính
các mức vi phạm $v_i(x^k)$ và các gradient trên tập active $\{\nabla F_i(x^k):i\in
K^{\cvx}_k\cup K^{\ps}_k\}$, lập hai hướng tìm kiếm $g^{\cvx}_k,g^{\ps}_k$ dưới dạng tổng có trọng số
của các gradient này, tính các cỡ bước, và thực hiện một phép chiếu lên $K$.
Thuật toán không dùng ước lượng chuẩn toán tử, cũng không dùng tìm kiếm theo tia. Chi phí chủ yếu nằm ở
việc tính các ràng buộc và phép chiếu. Chi phí cho bài toán chọn đặc trưng
được nêu ở Mục~\ref{sec:fs}
(Nhận xét~\ref{rem:fs-cost}).

\subsection{Các bổ đề bổ trợ}\label{ssec:aux}
Hai nhánh cần những lập luận khác nhau. Trên nhánh lồi, Bổ đề~\ref{lem:grad-cvx} chặn dưới
$\inner{\nabla F_i(x)}{x-x^\star}$ bởi mức vi phạm $v_i(x)$; nhờ đó $\norm{x^k-x^\star}^2$ giảm
một lượng cỡ $\mu_kT_k$. Trên nhánh giả lồi,
định nghĩa~\eqref{eq:contrapositive} chỉ cho biết dấu của tích vô hướng này.
Bổ đề~\ref{lem:fejer} gộp cả hai nhánh vào một bất đẳng thức tựa Fej\'er; cách chọn
tham số $\zeta$ trong bất đẳng thức Young giữ cho số hạng của nhánh lồi luôn âm. Sau đó Giai đoạn~2 của
Định lý~\ref{thm:strong} chứng minh bằng phản chứng rằng có ít nhất một điểm tụ
chấp nhận được, và Giai đoạn~3 chứng minh toàn bộ dãy hội tụ mạnh tới điểm đó. Hai trường hợp suy biến
được suy ra dưới dạng hệ quả.

Trong toàn bộ phần sau, $\{x^k\}$ là dãy vô hạn sinh bởi Thuật toán~\ref{alg:dtap} khi
bỏ phép kiểm tra dừng. Như được chứng minh dưới đây
(Bổ đề~\ref{lem:jcvx0} và Định lý~\ref{thm:strong}), $J(x^k)\to 0$; do đó với mọi
$\varepsilon>0$ tồn tại chỉ số hữu hạn $k$ với $J(x^k)\le\varepsilon$, nên phiên bản có phép kiểm tra dừng
kết thúc sau hữu hạn bước với một điểm $\varepsilon$-chấp nhận được.

Ta bắt đầu với hai kết quả chuẩn bị được dùng nhiều lần ở phần sau.

\begin{lemma}\label{lem:proj-nonexp}
Cho $K\subseteq H$ là tập lồi đóng khác rỗng. Với mọi $y\in H$ và $z\in K$,
\begin{equation}\label{eq:proj-nonexp}
\norm{P_K(y)-z}^2\le\norm{y-z}^2-\norm{y-P_K(y)}^2\le\norm{y-z}^2 .
\end{equation}
\end{lemma}
\begin{proof}
Đặt $p:=P_K(y)$. Theo Mệnh đề~\ref{prop:proj-char}, $\inner{y-p}{z-p}\le0$, nên khai triển
$\norm{y-z}^2=\norm{y-p}^2+2\inner{y-p}{p-z}+\norm{p-z}^2$ và dùng $\inner{y-p}{p-z}\ge0$
ta được $\norm{y-z}^2\ge\norm{y-p}^2+\norm{p-z}^2$, chính là~\eqref{eq:proj-nonexp}.
\end{proof}

\begin{lemma}\label{lem:seq}
Cho $\{a_k\}\subset[0,\infty)$ và $\{\varepsilon_k\}\subset[0,\infty)$ thỏa
$a_{k+1}\le a_k+\varepsilon_k$ với mọi $k$ và $\sum_{k=0}^\infty\varepsilon_k<\infty$.
Khi đó $\{a_k\}$ hội tụ tới một giới hạn hữu hạn.
\end{lemma}
\begin{proof}
Dãy $b_k:=a_k-\sum_{l<k}\varepsilon_l$ không tăng (do $a_{k+1}\le a_k+\varepsilon_k$) và
bị chặn dưới bởi $-\sum_l\varepsilon_l>-\infty$, nên hội tụ; vì
$\sum_{l<k}\varepsilon_l\to S<\infty$, dãy $a_k=b_k+\sum_{l<k}\varepsilon_l$ hội tụ.
\end{proof}

Trên nhánh lồi, tính lồi cho một liên hệ định lượng giữa gradient và
mức vi phạm.

\begin{lemma}\label{lem:grad-cvx}
Với $x^\star\in\Omega$, $i\in I^{\cvx}$ và $x\in K$,
\[
v_i(x)\,\inner{\nabla F_i(x)}{x-x^\star}\ge v_i(x)^2 .
\]
\end{lemma}
\begin{proof}
Nếu $v_i(x)=0$ thì cả hai vế bằng $0$. Giả sử $v_i(x)>0$, tức là $F_i(x)>b_i$. Định nghĩa
$\varphi:[0,1]\to\R$ bởi $\varphi(t):=F_i((1-t)x+tx^\star)$; vì $K$ lồi,
$[x,x^\star]\subset K$, nên $\varphi$ xác định, theo (A1)(b) nó lồi, và
theo (A2) nó khả vi với $\varphi'(0^+)=\inner{\nabla F_i(x)}{x^\star-x}$. Bất đẳng thức
gradient một chiều cho $\varphi(1)\ge\varphi(0)+\varphi'(0^+)$, tức là
$F_i(x^\star)\ge F_i(x)+\inner{\nabla F_i(x)}{x^\star-x}$. Vì $x^\star\in\Omega$ nên
$F_i(x^\star)\le b_i<F_i(x)$, do đó
$\inner{\nabla F_i(x)}{x-x^\star}\ge F_i(x)-b_i=v_i(x)$. Nhân với $v_i(x)>0$ ta được
khẳng định.
\end{proof}

\begin{remark}\label{rem:grad-cvx-sum}
Cộng Bổ đề~\ref{lem:grad-cvx} trên tập active xấp xỉ $K^{\cvx}_k$ ta được
$\inner{g^{\cvx}_k}{x^k-x^\star}\ge\sum_{i\in K^{\cvx}_k}v_i(x^k)^2=2\Phi^{\cvx}_k$, một cận
được dùng trong chứng minh các Bổ đề~\ref{lem:fejer} và~\ref{lem:jcvx0}.
\end{remark}

Trên nhánh giả lồi, ta chỉ có thông tin về dấu.

\begin{lemma}\label{lem:sign-ps}
Với $x^\star\in\Omega$, $j\in I^{\ps}$ và $x\in K$,
$v_j(x)\,\inner{\nabla F_j(x)}{x-x^\star}\ge 0$, với bất đẳng thức chặt khi $v_j(x)>0$.
\end{lemma}
\begin{proof}
Nếu $v_j(x)=0$ thì cả hai vế bằng $0$. Nếu $v_j(x)>0$ thì $F_j(x)>b_j\ge F_j(x^\star)$. Theo
(A1)(b), $F_j|_K$ giả lồi; áp dụng~\eqref{eq:contrapositive}
với $y=x^\star$ ta được $\inner{\nabla F_j(x)}{x^\star-x}<0$, tức là
$\inner{\nabla F_j(x)}{x-x^\star}>0$. Nhân với $v_j(x)>0$ ta được bất đẳng thức
chặt.
\end{proof}

Cuối cùng là hai cận chuẩn của các hướng tìm kiếm, dựa trên tính compact.

\begin{lemma}\label{lem:bound-cvx}
Giả sử $I^{\cvx}\ne\emptyset$ và $B\subset U$ compact. Với
$L_B:=\max_{i\in I^{\cvx}}\sup_{x\in B}\norm{\nabla F_i(x)}$ (hữu hạn theo (A2)), với mọi
$x\in B$,
\[
\norm{g^{\cvx}(x)}^2\le 2\,\abs{I^{\cvx}}\,L_B^2\,\Phi^{\cvx}(x).
\]
\end{lemma}
\begin{proof}
Đặt $K'=K^{\cvx}(x)$, $m'=\abs{K'}\le\abs{I^{\cvx}}$. Theo bất đẳng thức tam giác và
Cauchy--Schwarz,
\[
\norm{g^{\cvx}(x)}=\Bigl\|\sum_{i\in K'}v_i(x)\nabla F_i(x)\Bigr\|
\le\Bigl(\sum_{i\in K'}v_i(x)^2\Bigr)^{1/2}\Bigl(\sum_{i\in K'}\norm{\nabla F_i(x)}^2\Bigr)^{1/2}.
\]
Vì $\sum_{i\in K'}v_i(x)^2=2\Phi^{\cvx}(x)$ và
$\sum_{i\in K'}\norm{\nabla F_i(x)}^2\le m'L_B^2\le\abs{I^{\cvx}}L_B^2$, bình phương lên ta được
khẳng định.
\end{proof}

\begin{lemma}\label{lem:bound-ps}
Giả sử $I^{\ps}\ne\emptyset$ và $B\subset U$ compact. Với
$L'_B:=\max_{j\in I^{\ps}}\sup_{x\in B}\norm{\nabla F_j(x)}$ và
$V_B:=\max_{j\in I^{\ps}}\sup_{x\in B}v_j(x)$ (cả hai hữu hạn theo (A2)), với mọi $x\in B$,
\[
\norm{g^{\ps}(x)}\le\abs{I^{\ps}}\,V_B\,L'_B .
\]
\end{lemma}
\begin{proof}
Theo bất đẳng thức tam giác,
$\norm{g^{\ps}(x)}=\bigl\|\sum_{j\in K^{\ps}(x)}v_j(x)\nabla F_j(x)\bigr\|
\le\sum_{j\in K^{\ps}(x)}v_j(x)\norm{\nabla F_j(x)}\le\abs{K^{\ps}(x)}V_BL'_B
\le\abs{I^{\ps}}V_BL'_B$.
\end{proof}

\subsection{Tính bị chặn và sự giảm của sai số lặp}\label{ssec:bounded}
Bất đẳng thức sau đây thuộc loại tựa Fej\'er theo nghĩa của
Combettes~\citep{combettes2001quasifejer}. Nó cho tính bị chặn của $\{x^k\}$ và sự giảm của
mức vi phạm trên nhánh lồi. Đặt
\begin{equation}\label{eq:c0C0}
c_0:=\frac{\rho_{\cvx}(2-\rho_{\cvx})}{2}\in\Bigl(0,\tfrac12\Bigr],\qquad
C_0:=\frac{2+\rho_{\cvx}}{2-\rho_{\cvx}}\in(1,\infty),
\end{equation}
và
\begin{equation}\label{eq:Tk}
T_k:=\begin{cases}\dfrac{(2\Phi^{\cvx}_k)^2}{\norm{g^{\cvx}_k}^2} & \text{nếu }\norm{g^{\cvx}_k}>0,\\ 0 & \text{nếu }\norm{g^{\cvx}_k}=0.\end{cases}
\end{equation}
Với $x^\star\in\Omega$ cố định, viết $R^{x^\star}_k:=\inner{g^{\ps}_k}{x^k-x^\star}$.

\begin{lemma}\label{lem:fejer}
Với mọi $x^\star\in\Omega$ và $k\ge 0$,
\begin{equation}\label{eq:fejer}
\norm{x^{k+1}-x^\star}^2\le\norm{x^k-x^\star}^2-c_0\mu_k T_k-2(1-\mu_k)\eta^{\ps}_k R^{x^\star}_k+C_0\gamma_k^2 .
\end{equation}
\end{lemma}
\begin{proof}
Áp dụng Bổ đề~\ref{lem:proj-nonexp} với $y=y^k$ và $z=x^\star\in K$, và dùng
$y^k=x^k-d_k$ từ~\eqref{eq:dk}:
\begin{equation}\label{eq:p1}
\norm{x^{k+1}-x^\star}^2\le\norm{x^k-x^\star}^2-2\inner{d_k}{x^k-x^\star}+\norm{d_k}^2 .
\end{equation}
\emph{Số hạng chéo.} Khai triển $d_k$,
\begin{equation}\label{eq:p2}
\inner{d_k}{x^k-x^\star}=\mu_k\eta^{\cvx}_k\inner{g^{\cvx}_k}{x^k-x^\star}+(1-\mu_k)\eta^{\ps}_k R^{x^\star}_k .
\end{equation}
Nếu $\norm{g^{\cvx}_k}=0$ thì $\eta^{\cvx}_k=0$ và số hạng đầu bằng $0$. Nếu
$\norm{g^{\cvx}_k}>0$, Nhận xét~\ref{rem:grad-cvx-sum} cho
$\inner{g^{\cvx}_k}{x^k-x^\star}\ge 2\Phi^{\cvx}_k$; nhân với
$\mu_k\eta^{\cvx}_k\ge 0$ và thay
$\eta^{\cvx}_k=2\rho_{\cvx}\Phi^{\cvx}_k/\norm{g^{\cvx}_k}^2$,
\begin{equation}\label{eq:p4}
\mu_k\eta^{\cvx}_k\inner{g^{\cvx}_k}{x^k-x^\star}\ge\mu_k\rho_{\cvx}\frac{(2\Phi^{\cvx}_k)^2}{\norm{g^{\cvx}_k}^2}=\mu_k\rho_{\cvx}T_k .
\end{equation}
\emph{Số hạng bậc hai.} Theo bất đẳng thức Young
$\norm{u+v}^2\le(1+\zeta)\norm{u}^2+(1+\zeta^{-1})\norm{v}^2$ với
$u=\mu_k\eta^{\cvx}_k g^{\cvx}_k$, $v=(1-\mu_k)\eta^{\ps}_k g^{\ps}_k$,
\[
\norm{d_k}^2\le(1+\zeta)\rho_{\cvx}^2\mu_k^2 T_k+(1+\zeta^{-1})(1-\mu_k)^2(\eta^{\ps}_k)^2\norm{g^{\ps}_k}^2 .
\]
Theo định nghĩa $\eta^{\ps}_k=\gamma_k/\max\{1,\norm{g^{\ps}_k}\}$, nên
$\eta^{\ps}_k\norm{g^{\ps}_k}\le\gamma_k$ và do đó
$(\eta^{\ps}_k)^2\norm{g^{\ps}_k}^2\le\gamma_k^2$; cùng với $(1-\mu_k)^2\le 1$,
\begin{equation}\label{eq:p6}
\norm{d_k}^2\le(1+\zeta)\rho_{\cvx}^2\mu_k^2 T_k+(1+\zeta^{-1})\gamma_k^2 .
\end{equation}
Chọn $\zeta:=(2-\rho_{\cvx})/(2\rho_{\cvx})>0$. Khi đó
$(1+\zeta)\rho_{\cvx}=(2+\rho_{\cvx})/2$, nên
$(1+\zeta)\rho_{\cvx}^2=\rho_{\cvx}(2+\rho_{\cvx})/2$, và
$1+\zeta^{-1}=(2+\rho_{\cvx})/(2-\rho_{\cvx})=C_0$. Kết hợp
\eqref{eq:p1}, \eqref{eq:p2}, \eqref{eq:p4} và~\eqref{eq:p6},
\[
\norm{x^{k+1}-x^\star}^2\le\norm{x^k-x^\star}^2+\bigl[-2\mu_k\rho_{\cvx}+(1+\zeta)\rho_{\cvx}^2\mu_k^2\bigr]T_k-2(1-\mu_k)\eta^{\ps}_k R^{x^\star}_k+C_0\gamma_k^2 .
\]
Biểu thức trong ngoặc vuông bằng
$\mu_k\rho_{\cvx}\bigl(-2+\tfrac{2+\rho_{\cvx}}{2}\mu_k\bigr)$. Vì $\rho_{\cvx}<2$, với
mọi $\mu_k\in[0,1]$ thừa số $-2+\tfrac{2+\rho_{\cvx}}{2}\mu_k$ không vượt quá
$-(2-\rho_{\cvx})/2<0$ (giá trị lớn nhất trên $[0,1]$ đạt tại $\mu_k=1$); nhân với
$\mu_k\rho_{\cvx}\ge0$ ta được biểu thức trong ngoặc $\le -c_0\mu_k$. Từ đó suy ra~\eqref{eq:fejer}.
\end{proof}

\begin{lemma}\label{lem:bounded}
Dãy $\{x^k\}$ bị chặn. Hơn nữa, với mỗi $x^\star\in\Omega$ cố định:
\textup{(i)} $\{\norm{x^k-x^\star}^2\}$ hội tụ trong $[0,\infty)$;
\textup{(ii)} $\sum_{k=0}^\infty\mu_k T_k<\infty$;
\textup{(iii)} $\sum_{k=0}^\infty(1-\mu_k)\eta^{\ps}_k R^{x^\star}_k<\infty$.
\end{lemma}
\begin{proof}
Cố định $x^\star\in\Omega$ và đặt $a_k:=\norm{x^k-x^\star}^2$. Theo Bổ đề~\ref{lem:sign-ps},
$R^{x^\star}_k\ge0$, và $\mu_k,T_k,(1-\mu_k)\eta^{\ps}_k\ge0$, nên~\eqref{eq:fejer} cho
$a_{k+1}\le a_k+C_0\gamma_k^2$. Vì $\sum_k C_0\gamma_k^2<\infty$ theo (A3)(iii),
Bổ đề~\ref{lem:seq} cho (i), và suy ra $\{x^k\}$ bị chặn. Cộng
\eqref{eq:fejer} từ $0$ tới $N-1$,
\[
a_N+\sum_{k=0}^{N-1}\bigl(c_0\mu_k T_k+2(1-\mu_k)\eta^{\ps}_k R^{x^\star}_k\bigr)\le a_0+C_0\sum_{k=0}^{N-1}\gamma_k^2 .
\]
Vế phải bị chặn đều theo $N$; cho $N\to\infty$ và dùng $c_0>0$ cùng với
các số hạng không âm ta được (ii) và (iii).
\end{proof}

\begin{lemma}\label{lem:jcvx0}
$\displaystyle\lim_{k\to\infty}J^{\cvx}_k=0$.
\end{lemma}
\begin{proof}
Nếu $I^{\cvx}=\emptyset$ thì $J^{\cvx}_k=0$ với mọi $k$. Giả sử $I^{\cvx}\ne\emptyset$, cố định
$x^\star\in\Omega$, và đặt $B:=\overline{\conv}\{x^k:k\ge0\}$, compact theo Bổ đề~\ref{lem:bounded} và nằm
trong $K\subset U$. Theo (A2) và
tính compact,
\[
L:=\max_{i\in I^{\cvx}}\sup_{x\in B}\norm{\nabla F_i(x)}<\infty,\qquad
V:=\max_{i\in I}\sup_{x\in B}\abs{F_i(x)-b_i}<\infty.
\]
Nếu $V=0$ thì mọi $v_i(x^k)=0$ và
$J^{\cvx}_k=0$; giả sử $V>0$. Nếu $L=0$ thì mỗi $F_i$, $i\in I^{\cvx}$, là hằng trên
$B$, nên $g^{\cvx}_k=0$ và, theo Nhận xét~\ref{rem:grad-cvx-sum},
$0=\inner{g^{\cvx}_k}{x^k-x^\star}\ge 2\Phi^{\cvx}_k\ge(J^{\cvx}_k)^2$, cho ta
$J^{\cvx}_k=0$; giả sử $L>0$. Đặt $m^{\cvx}:=\abs{I^{\cvx}}\ge1$.

Xét một bước với $J^{\cvx}_k>0$ và chọn $i^\star\in I^{\cvx}$ với
$v_{i^\star}(x^k)=J^{\cvx}_k$. Vì $v_{i^\star}(x^k)\ge J^{\cvx}_k-\delta$, ta có
$i^\star\in K^{\cvx}_k$, do đó
\begin{equation}\label{eq:p8}
\Phi^{\cvx}_k\ge\tfrac12 v_{i^\star}(x^k)^2=\tfrac12(J^{\cvx}_k)^2>0 .
\end{equation}
Nhận xét~\ref{rem:grad-cvx-sum} cho
$\inner{g^{\cvx}_k}{x^k-x^\star}\ge2\Phi^{\cvx}_k>0$, nên $g^{\cvx}_k\ne0$; cùng với
Bổ đề~\ref{lem:bound-cvx}, $\norm{g^{\cvx}_k}^2\le2m^{\cvx}L^2\Phi^{\cvx}_k$, và cùng với
\eqref{eq:p8},
\begin{equation}\label{eq:p9}
T_k=\frac{(2\Phi^{\cvx}_k)^2}{\norm{g^{\cvx}_k}^2}\ge\frac{2\Phi^{\cvx}_k}{m^{\cvx}L^2}\ge\frac{(J^{\cvx}_k)^2}{m^{\cvx}L^2}.
\end{equation}
Để chặn dưới $\mu_k$: với mọi $i\in I$ và $x^k\in B$,
$v_i(x^k)\le\abs{F_i(x^k)-b_i}\le V$, nên $J^{\cvx}_k\le V$ và $J^{\ps}_k\le V$; khi
$J^{\cvx}_k>0$,
\begin{equation}\label{eq:p10}
\mu_k=\frac{J^{\cvx}_k}{J^{\cvx}_k+J^{\ps}_k}\ge\frac{J^{\cvx}_k}{2V}.
\end{equation}
Kết hợp~\eqref{eq:p9} và~\eqref{eq:p10}, khi $J^{\cvx}_k>0$,
$\mu_k T_k\ge(J^{\cvx}_k)^3/(2Vm^{\cvx}L^2)$; bất đẳng thức này cũng đúng (dạng $0\ge0$) khi
$J^{\cvx}_k=0$. Theo Bổ đề~\ref{lem:bounded}(ii),
\begin{equation}\label{eq:p12}
\sum_{k=0}^\infty(J^{\cvx}_k)^3\le 2Vm^{\cvx}L^2\sum_{k=0}^\infty\mu_k T_k<\infty,
\end{equation}
nên $(J^{\cvx}_k)^3\to0$ và do đó $J^{\cvx}_k\to0$.
\end{proof}

\begin{proposition}\label{prop:regularity}
\[
\sum_{k=0}^{\infty}\norm{x^{k+1}-x^k}^{2}<\infty;
\]
nói riêng,
$\norm{x^{k+1}-x^k}\to0$.
\end{proposition}
\begin{proof}
Vì $x^k\in K$ và $x^{k+1}=P_K(y^k)$, Bổ đề~\ref{lem:proj-nonexp} với $y=y^k$, $z=x^k$
cho $\norm{x^{k+1}-x^k}\le\norm{y^k-x^k}=\norm{d_k}$. Theo~\eqref{eq:p6} và
$\mu_k^2\le\mu_k$, $\norm{d_k}^2\le(1+\zeta)\rho_{\cvx}^2\mu_k T_k+C_0\gamma_k^2$.
Lấy tổng và dùng Bổ đề~\ref{lem:bounded}(ii) cùng (A3)(iii), cả hai chuỗi đều hữu hạn.
\end{proof}

\subsection{Hội tụ mạnh tới một điểm chấp nhận được}\label{ssec:strong}
Bây giờ ta chứng minh toàn bộ dãy hội tụ mạnh tới một điểm chấp nhận được. Chứng minh có ba
giai đoạn. Giai đoạn~1 dùng Bổ đề~\ref{lem:jcvx0} để chứng minh mọi điểm tụ thỏa các
ràng buộc lồi. Giai đoạn~2 chứng minh, bằng phản chứng và tính giả lồi, rằng có ít nhất một
điểm tụ chấp nhận được. Giai đoạn~3 dùng sự hội tụ của $\norm{x^k-\bar x}$ để chứng minh
toàn bộ dãy hội tụ tới điểm đó.

\begin{theorem}\label{thm:strong}
Dưới Giả thiết~\ref{ass:standing}, dãy $\{x^k\}$ sinh bởi
Thuật toán~\ref{alg:dtap} hội tụ mạnh tới một điểm $x^\infty\in\Omega$.
\end{theorem}
\begin{proof}
Theo Bổ đề~\ref{lem:bounded}, $\{x^k\}$ bị chặn; gọi
$C:=\{\bar x\in H:\exists\,k_l\to\infty,\ x^{k_l}\to\bar x\}$ là tập điểm tụ (khác rỗng, compact)
của nó. Vì $K$ đóng và $x^k\in K$, ta có $C\subseteq K$.

\emph{Giai đoạn 1 (các ràng buộc lồi).} Cố định $\bar x\in C$ và $x^{k_l}\to\bar x$. Theo (A2) mỗi
$F_i$ liên tục, nên $v_i(x^{k_l})\to v_i(\bar x)$; với $i\in I^{\cvx}$,
$v_i(x^{k_l})\le J^{\cvx}_{k_l}\to0$ theo Bổ đề~\ref{lem:jcvx0}, nên theo nguyên lý kẹp
$v_i(\bar x)=0$, tức là $F_i(\bar x)\le b_i$. Do đó
\begin{equation}\label{eq:p13}
C\subseteq\bigcap_{i\in I^{\cvx}}\{x\in K:F_i(x)\le b_i\}.
\end{equation}
Nếu $I^{\ps}=\emptyset$ thì~\eqref{eq:p13} cho $C\subseteq\Omega$,
nên $C\cap\Omega=C\ne\emptyset$.
Từ đây giả sử $I^{\ps}\ne\emptyset$.

\emph{Giai đoạn 2 (một điểm tụ chấp nhận được).} Giả sử ngược lại rằng
$C\cap\Omega=\emptyset$. Theo~\eqref{eq:p13} và phép chia $I=I^{\cvx}\cup I^{\ps}$, với
mỗi $\bar x\in C$ có $j\in I^{\ps}$ với $F_j(\bar x)>b_j$, tức là
$J^{\ps}(\bar x)>0$. Vì là giá trị lớn nhất của hữu hạn hàm liên tục, $J^{\ps}$
liên tục và dương trên tập compact $C$, nên $\alpha:=\min_{x\in C}J^{\ps}(x)>0$.

\smallskip
\emph{Khẳng định 1 (cận dưới đều cho $J^{\ps}_k$).} Tồn tại $\tau_0>0$ và $k_0$
sao cho $J^{\ps}_k\ge\tau_0$ với mọi $k\ge k_0$. Thật vậy, đặt $\tau_0:=\alpha/2$. Nếu khẳng định
sai thì có $k_l\to\infty$ với $J^{\ps}_{k_l}<\tau_0$; trích dãy con (giữ nguyên
ký hiệu) $x^{k_l}\to\tilde x\in C$, tính liên tục cho $J^{\ps}(\tilde x)\le\tau_0<\alpha$,
mâu thuẫn với định nghĩa của $\alpha$.

\smallskip
\emph{Khẳng định 2 ($\mu_k\to0$).} Với $k\ge k_0$,
$\mu_k=J^{\cvx}_k/(J^{\cvx}_k+J^{\ps}_k)\le J^{\cvx}_k/\tau_0\to0$ theo
Bổ đề~\ref{lem:jcvx0}. Do đó có $k_1\ge k_0$ với $1-\mu_k\ge1/2$ với mọi $k\ge k_1$.

\smallskip
\emph{Khẳng định 3 (cận dưới đều cho $R^{x^\star}_k$).} Cố định $x^\star\in\Omega$. Tồn tại
$\xi_0>0$ và $k_2\ge k_1$ với $R^{x^\star}_k\ge\xi_0$ với mọi $k\ge k_2$. Giả sử
ngược lại: tồn tại $k_l\to\infty$ với $R^{x^\star}_{k_l}\to0$. Vì $\{x^{k_l}\}$ bị chặn,
trích $x^{k_l}\to\bar x\in C$; vì $I^{\ps}$ chỉ có hữu hạn tập con, trích tiếp
để $K^{\ps}_{k_l}\equiv K^\star$ không đổi. Theo giả thiết phản chứng
$J^{\ps}(\bar x)>0$; chọn $j^\star\in\argmax_{j\in I^{\ps}}v_j(\bar x)$, nên
$v_{j^\star}(\bar x)=J^{\ps}(\bar x)>0$. Theo tính liên tục $v_{j^\star}(x^{k_l})\to
v_{j^\star}(\bar x)$ và $J^{\ps}_{k_l}\to J^{\ps}(\bar x)$, nên
$J^{\ps}_{k_l}-v_{j^\star}(x^{k_l})\to0$; với $\delta>0$ từ (A3)(i), khi $l$ đủ lớn ta
có $v_{j^\star}(x^{k_l})\ge J^{\ps}_{k_l}-\delta$, do đó $j^\star\in K^{\ps}_{k_l}=K^\star$.
Mỗi số hạng của
$R^{x^\star}_{k_l}=\sum_{j\in K^\star}v_j(x^{k_l})\inner{\nabla F_j(x^{k_l})}{x^{k_l}-x^\star}$
không âm theo Bổ đề~\ref{lem:sign-ps}, nên chỉ giữ lại $j^\star$,
\[
R^{x^\star}_{k_l}\ge v_{j^\star}(x^{k_l})\inner{\nabla F_{j^\star}(x^{k_l})}{x^{k_l}-x^\star}.
\]
Cho $l\to\infty$ và dùng tính liên tục của $v_{j^\star}$ và $\nabla F_{j^\star}$,
\[
\liminf_{l\to\infty}R^{x^\star}_{k_l}\ge v_{j^\star}(\bar x)\inner{\nabla F_{j^\star}(\bar x)}{\bar x-x^\star}=:\Lambda .
\]
Vì $j^\star\in I^{\ps}$ và $x^\star\in\Omega$, $F_{j^\star}(x^\star)\le b_{j^\star}<F_{j^\star}(\bar x)$;
định nghĩa~\eqref{eq:contrapositive} với $x=\bar x$, $y=x^\star$ cho
$\inner{\nabla F_{j^\star}(\bar x)}{\bar x-x^\star}>0$, và cùng với $v_{j^\star}(\bar x)>0$ ta
được $\Lambda>0$, mâu thuẫn với $R^{x^\star}_{k_l}\to0$. Khẳng định~3 được chứng minh.

\smallskip
\emph{Đi đến mâu thuẫn.} Theo Bổ đề~\ref{lem:bound-ps} trên
$B:=\overline{\conv}\{x^k\}$, $\norm{g^{\ps}_k}\le\abs{I^{\ps}}V_B L'_B$ với
$L'_B:=\max_{j\in I^{\ps}}\sup_B\norm{\nabla F_j}<\infty$. Đặt
$G':=\max\{1,\abs{I^{\ps}}V_B L'_B\}$; khi đó $\max\{1,\norm{g^{\ps}_k}\}\le G'$ với mọi $k$,
nên $\eta^{\ps}_k=\gamma_k/\max\{1,\norm{g^{\ps}_k}\}\ge\gamma_k/G'$. Với $k\ge k_2$,
kết hợp với $1-\mu_k\ge1/2$ và $R^{x^\star}_k\ge\xi_0$,
\[
\sum_{k=0}^\infty(1-\mu_k)\eta^{\ps}_k R^{x^\star}_k\ge\sum_{k\ge k_2}\frac12\cdot\frac{\gamma_k}{G'}\cdot\xi_0=\frac{\xi_0}{2G'}\sum_{k\ge k_2}\gamma_k=+\infty
\]
theo (A3)(iii), mâu thuẫn với Bổ đề~\ref{lem:bounded}(iii). Do đó $C\cap\Omega\ne\emptyset$.

\emph{Giai đoạn 3 (hội tụ của toàn bộ dãy).} Cố định $\bar x\in C\cap\Omega$.
Bổ đề~\ref{lem:bounded}(i), áp dụng với $x^\star=\bar x$, bảo đảm
$\lim_{k\to\infty}\norm{x^k-\bar x}^2$ tồn tại. Vì $\bar x\in C$, có một dãy con
$x^{k_l}$ hội tụ tới $\bar x$. Do đó
\[
\lim_{k\to\infty}\norm{x^k-\bar x}^2
=\lim_{l\to\infty}\norm{x^{k_l}-\bar x}^2=0.
\]
Vậy toàn bộ dãy hội tụ tới $\bar x$, và ta đặt $x^\infty:=\bar x\in\Omega$.
Lập luận không cần $\Omega$ chỉ có một điểm: một điểm tụ chấp nhận được và
sự hội tụ của khoảng cách tới nó là đủ.
\end{proof}

\begin{remark}\label{rem:delta-role}
Tham số $\delta>0$ không cần thiết cho sự hội tụ. Khi $\delta=0$, ta lập luận theo
dãy con: dọc theo $x^{k_l}\to\bar x\in C$, chọn
$j_l\in\argmax_{j\in I^{\ps}}v_j(x^{k_l})$; vì $I^{\ps}$ hữu hạn, trích
$j_l\equiv j^\star$, và $v_{j^\star}(x^{k_l})=J^{\ps}(x^{k_l})\to J^{\ps}(\bar x)>0$, nên
$j^\star$ đạt mức vi phạm lớn nhất tại $\bar x$ và lập luận giả lồi
vẫn đúng. Vai trò của $\delta>0$ là làm chứng minh đơn giản hơn: một chỉ số đạt max tại $\bar x$ tự động
nằm trong $K^{\ps}_{k_l}$ khi $l$ đủ lớn, nên không cần bước trích dãy con dựa trên tính hữu hạn của tập chỉ số.
$\delta$ lớn hơn làm rộng $K^{\cvx}_k$ và $K^{\ps}_k$, điều này có thể cải thiện độ ổn định
số khi nhiều $v_i(x^k)$ gần bằng nhau.
\end{remark}

\begin{remark}\label{rem:rate}
Từ~\eqref{eq:p12}, $\sum_k(J^{\cvx}_k)^3<\infty$, suy ra
$\min_{0\le k<N}J^{\cvx}_k=O(N^{-1/3})$. Đây là cận cho mức vi phạm \emph{tốt nhất} của nhánh lồi
trong $N$ điểm lặp đầu tiên. Nó không phải là cận cho điểm lặp cuối, không phải tốc độ cho toàn bộ phần dư
$J(x^k)$, cũng không phải tốc độ cho khoảng cách tới một điểm chấp nhận được. Với các đại lượng này,
chỉ với (A1)--(A3), chúng tôi chỉ khẳng định sự hội tụ định tính.
\end{remark}

\subsection{Các cấu trúc ràng buộc suy biến}\label{ssec:degenerate}
Khi cấu trúc ràng buộc suy biến (chỉ có nhánh lồi hoặc chỉ có nhánh giả lồi),
\dtap{} và phân tích của nó trở nên đơn giản hơn, và ta được hai hệ quả của
Định lý~\ref{thm:strong}.

\begin{corollary}\label{cor:cvx}
Nếu $I^{\ps}=\emptyset$ thì $J^{\ps}_k\equiv0$, $g^{\ps}_k\equiv0$, và $\mu_k\equiv1$
mỗi khi $J^{\cvx}_k>0$, nên Thuật toán~\ref{alg:dtap} trở thành
$x^{k+1}=P_K(x^k-\eta^{\cvx}_k g^{\cvx}_k)$. Dãy hội tụ mạnh tới một điểm
$x^\infty\in\Omega=\{x\in K:F_i(x)\le b_i\text{ với mọi }i\in I^{\cvx}\}$.
\end{corollary}
\begin{proof}
Ở đây $d_k=\eta^{\cvx}_k g^{\cvx}_k$, và~\eqref{eq:p1} cho trực tiếp
$\norm{x^{k+1}-x^\star}^2\le\norm{x^k-x^\star}^2-2\eta^{\cvx}_k\inner{g^{\cvx}_k}{x^k-x^\star}+(\eta^{\cvx}_k)^2\norm{g^{\cvx}_k}^2$.
Nếu $\norm{g^{\cvx}_k}=0$ thì $\eta^{\cvx}_k=0$, $T_k=0$, và cận trở thành
$\norm{x^{k+1}-x^\star}^2\le\norm{x^k-x^\star}^2$. Ngược lại, Bổ đề~\ref{lem:grad-cvx} cho
$\inner{g^{\cvx}_k}{x^k-x^\star}\ge2\Phi^{\cvx}_k$, và thay
$\eta^{\cvx}_k=2\rho_{\cvx}\Phi^{\cvx}_k/\norm{g^{\cvx}_k}^2$ cùng với
$(\eta^{\cvx}_k)^2\norm{g^{\cvx}_k}^2=\eta^{\cvx}_k\cdot2\rho_{\cvx}\Phi^{\cvx}_k$ ta được, sau khi
rút gọn,
$2\eta^{\cvx}_k\inner{g^{\cvx}_k}{x^k-x^\star}-(\eta^{\cvx}_k)^2\norm{g^{\cvx}_k}^2\ge\rho_{\cvx}(2-\rho_{\cvx})(2\Phi^{\cvx}_k)^2/\norm{g^{\cvx}_k}^2=2c_0T_k$.
Do đó
\begin{equation}\label{eq:p15}
\norm{x^{k+1}-x^\star}^2\le\norm{x^k-x^\star}^2-2c_0T_k ,
\end{equation}
đúng với mọi $k$. Như vậy $\{\norm{x^k-x^\star}^2\}$ không tăng và hội tụ, và
lấy tổng ta được $\sum_k T_k<\infty$. Như trong Bổ đề~\ref{lem:jcvx0}, khi $J^{\cvx}_k>0$ ta có
$T_k\ge(J^{\cvx}_k)^2/(m^{\cvx}L^2)$ (để ý số mũ $2$ ở đây, so với $3$
trong~\eqref{eq:p12}, vì không có thừa số $\mu_k$), nên $\sum_k(J^{\cvx}_k)^2<\infty$ và
$J^{\cvx}_k\to0$. Cuối cùng, $I^{\ps}=\emptyset$ khiến Giai đoạn~2 của
Định lý~\ref{thm:strong} không còn cần thiết (thật vậy $C\subseteq\Omega$ đã có
từ~\eqref{eq:p13}), và Giai đoạn~3 cho $x^k\to x^\infty\in\Omega$.
\end{proof}

\begin{corollary}\label{cor:ps}
Nếu $I^{\cvx}=\emptyset$ thì $J^{\cvx}_k\equiv0$, $g^{\cvx}_k\equiv0$, và $\mu_k\equiv0$
mỗi khi $J^{\ps}_k>0$, nên Thuật toán~\ref{alg:dtap} trở thành
$x^{k+1}=P_K(x^k-\eta^{\ps}_k g^{\ps}_k)$ với $\eta^{\ps}_k=\gamma_k/\max\{1,\norm{g^{\ps}_k}\}$.
Dãy hội tụ mạnh tới một điểm $x^\infty\in\Omega$.
\end{corollary}
\begin{proof}
Lúc này $d_k=\eta^{\ps}_k g^{\ps}_k$ với $\norm{d_k}^2\le\gamma_k^2$, nên~\eqref{eq:p1} cho
$\norm{x^{k+1}-x^\star}^2\le\norm{x^k-x^\star}^2-2\eta^{\ps}_k R^{x^\star}_k+\gamma_k^2$,
chính là~\eqref{eq:fejer} trong trường hợp suy biến này ($T_k=0$, $1-\mu_k=1$). Vì
$R^{x^\star}_k\ge0$ (Bổ đề~\ref{lem:sign-ps}), $a_{k+1}\le a_k+\gamma_k^2$, và
Bổ đề~\ref{lem:seq} cho sự hội tụ của $\{\norm{x^k-x^\star}^2\}$ và tính bị chặn;
lấy tổng ta được $\sum_k\eta^{\ps}_k R^{x^\star}_k<\infty$. Giai đoạn~1 của
Định lý~\ref{thm:strong} cho $C\subseteq K$ (phần về ràng buộc lồi không có nội dung); Giai đoạn~2 với
$\mu_k=0$ (nên $1-\mu_k=1$) cho cùng mâu thuẫn $\sum\eta^{\ps}_k R^{x^\star}_k=+\infty$,
nên $C\cap\Omega\ne\emptyset$; Giai đoạn~3 cho $x^k\to x^\infty\in\Omega$.
\end{proof}


\section{Thực nghiệm số}\label{sec:experiments}
Chúng tôi tách riêng hiệu quả dự đoán và việc giải số hệ chấp nhận. Phép so sánh phân loại dùng tám bộ dữ liệu. Phần chẩn đoán bộ giải còn dùng thêm Movement Libras và các bài toán giả lồi không lồi có kiểm soát.

\subsection{Mô hình chọn đặc trưng}\label{sec:fs}
Ta áp dụng NSFP với hộp chuẩn cho bài toán chọn đặc trưng bằng cách kết hợp các tiêu chí từ mô hình quy hoạch phân thức của Wang và cộng sự~\citep{wang2021fractional} và mô hình tập thô mờ AFRSFCV~\citep{wang2026afrsfcv}. Cho $\mathcal{X}=\{x_1,\ldots,x_n\}$ là các mẫu huấn luyện, $A=\{a_1,\ldots,a_p\}$ là các đặc trưng, $X\in\R^{n\times p}$ là ma trận dữ liệu, và $y\in\{1,\ldots,r\}^n$ là nhãn lớp. Ở đây $n$ là số mẫu; chỉ số dưới đánh số đặc trưng, chỉ số trên đánh số vòng lặp. Các trọng số đặc trưng thuộc đơn hình lồi compact $K=\Delta_{p-1}=\{w\in\R^p:w\ge0,\ \mathbf1_p^{\top}w=1\}$. Mọi đại lượng có giám sát dưới đây được ước lượng từ dữ liệu huấn luyện.

\subsubsection{Tiêu chí phân thức và N-Fscore}\label{ssec:wang}\label{sssec:wang-inputs}\label{sssec:nfscore}
Ma trận tương đồng có giám sát $S\in\R^{p\times p}$ có các phần tử
\begin{equation}\label{eq:simMI}
s_{ij}=\max\left\{0,\frac{I(Z_i;Z_j;C)}{H(Z_i)+H(Z_j)}\right\},\qquad i,j=1,\ldots,p,
\end{equation}
trong đó $Z_i$ là đặc trưng $a_i$ đã rời rạc hóa, $C$ là biến lớp, và $I(Z_i;Z_j;C)=I(Z_i;C)+I(Z_j;C)-I((Z_i,Z_j);C)$ là thông tin tương tác của chúng. Entropy và thông tin tương hỗ được ước lượng từ histogram độ rộng đều với mười khoảng cho mỗi đặc trưng~\citep{brown2012}. Ta đặt $s_{ij}=0$ nếu mẫu số bằng không. Ma trận dư thừa là
\begin{equation}\label{eq:matrixQ}
Q_R=S+\delta_QI_p,\qquad \delta_Q\ge-\min\{0,\lambda_{\min}(S)\},
\end{equation}
nên $Q_R\succeq0$. Mã cài đặt dùng $\delta_Q=\max\{0,-\lambda_{\min}(S)\}+10^{-10}$. Cho $n_t$, $\bar a_{kt}$ và $\sigma_{kt}^2$ là kích thước lớp, trung bình và phương sai tổng thể của đặc trưng $a_k$ trong lớp $t$, và $\bar a_k$ là trung bình chung của nó. Điểm Fisher là
\begin{equation}\label{eq:rhoFisher}
\rho_k=\frac{\sum_{t=1}^r n_t(\bar a_{kt}-\bar a_k)^2}{\sum_{t=1}^r n_t\sigma_{kt}^2+\varepsilon_\rho},\qquad \varepsilon_\rho=10^{-12}.
\end{equation}
Các điểm này không âm. Số hạng hiệu chỉnh ở mẫu số tránh chia cho không nhưng không bảo đảm $\rho_k>0$, vì tử số có thể bằng không.

Bài toán quy hoạch phân thức nằm sau N-Fscore là
\begin{equation}\label{eq:wangOrig}
\min_{w\in K}\ \frac{w^{\top}Q_Rw}{\rho^{\top}w},
\end{equation}
với giả thiết mẫu số dương trên $K$. Tử số đo độ dư thừa và mẫu số đo độ liên quan. Phương pháp cơ sở N-Fscore của chúng tôi dùng giai đoạn mạng nơ-ron chiếu của~\citep{wang2021fractional}, với điểm Fisher trong~\eqref{eq:rhoFisher}. Sau đó nó áp dụng quy tắc chọn chung~\eqref{eq:thetaRule}, thay cho giai đoạn $k$-winners-take-all của phương pháp gốc. Động lực học và cách cài đặt Euler hiện của nó được nêu ở Mục~\ref{sssec:baselines}.

\subsubsection{Tiêu chí tập thô mờ và AFRSFCV}\label{ssec:afrsfcv}\label{sssec:afrsfcv-inputs}\label{sssec:afrsfcv-model}
Với mỗi lớp, viết $D_t=\{x_i:y_i=t\}$ và $D_t(x_j)=\mathbf1_{\{y_j=t\}}$. AFRSFCV đưa phương sai lớp vào một quan hệ mờ. Với $a_k(x_i)=X_{ik}$, quan hệ được cài đặt là
\begin{equation}\label{eq:simRel}
R_k(x_i,x_j)=\exp\!\left(-\frac{[a_k(x_i)-a_k(x_j)]^2}{2\widetilde\sigma_{k,y_i}^2}\right),
\qquad \widetilde\sigma_{kt}^2=\max\{\sigma_{kt}^2,10^{-10}\}.
\end{equation}
Ngưỡng dưới của phương sai bảo đảm công thức có nghĩa khi một đặc trưng là hằng trong một lớp. Xấp xỉ dưới mờ là
\begin{equation}\label{eq:lower}
\underline{R_k}(D_t)(x_i)=\min_{x_j\in\mathcal{X}}\max\{1-R_k(x_i,x_j),D_t(x_j)\}.
\end{equation}
Các giá trị này tạo thành ma trận $\underline{R_k}(D)\in[0,1]^{n\times r}$. Cho $Y_{it}=\mathbf1_{\{y_i=t\}}$ là ma trận nhãn và định nghĩa ma trận quyết định bởi
\begin{equation}\label{eq:Pmat}
P_{ik}=\mathbf1_{\{\argmax_{1\le t\le r}\underline{R_k}(D_t)(x_i)=y_i\}}.
\end{equation}
Khi nhiều lớp cùng đạt giá trị lớn nhất, ta chọn lớp có chỉ số nhỏ nhất. Như vậy $P_{ik}=1$ khi và chỉ khi đặc trưng $a_k$ dự đoán đúng mẫu $x_i$ theo quy tắc này.

Phương pháp cơ sở AFRSFCV~\citep{wang2026afrsfcv} cực tiểu hóa
\begin{equation}\label{eq:afrsfcvLoss}
\begin{split}
\min_{\tau\in\R^p}\ L(\tau)
&=\frac12\left\|Y-\sum_{k=1}^{p}w_k(\tau)\underline{R_k}(D)\right\|_F^2\\
&\quad+\frac{\lambda_1}{2}\|\mathbf1_n-Pw(\tau)\|^2+\frac{\lambda_2}{2}\|w(\tau)\|^2,
\end{split}
\end{equation}
trong đó $w_k(\tau)=e^{\tau_k}/\sum_{l=1}^{p}e^{\tau_l}$ bảo đảm $w(\tau)\in K$. Ba số hạng là sai số xấp xỉ nhãn, trọng số của các đặc trưng phân loại sai trên mỗi mẫu, và một số hạng Tikhonov. Phương pháp cơ sở dùng hạ gradient với số vòng lặp cố định, rồi áp dụng~\eqref{eq:thetaRule}. Gradient và bước cập nhật được nêu ở Mục~\ref{sssec:baselines}.

\subsubsection{Bài toán chấp nhận và chọn đặc trưng}\label{ssec:reduction}\label{ssec:selection}
Ta đặt các cận riêng rẽ lên ba tiêu chí,
\begin{align}
F_1(w)&=\left\|Y-\sum_{k=1}^{p}w_k\underline{R_k}(D)\right\|_F^2,\label{eq:fsF1}\\
F_2(w)&=\|\mathbf1_n-Pw\|^2,\label{eq:fsF2}\\
F_3(w)&=\frac{w^{\top}Q_Rw}{\rho^{\top}w},\qquad U_\rho=\{w:\rho^{\top}w>0\}.\label{eq:fsF3}
\end{align}
Ở đây $F_1,F_2$ là hai số hạng đầu của~\eqref{eq:afrsfcvLoss}, bỏ hệ số, và $F_3$ là hàm mục tiêu trong~\eqref{eq:wangOrig}. Với các ngưỡng $b_i>0$, bài toán thu được là
\begin{equation}\label{eq:nsfp-fs}
\text{tìm }w^\star\in\Omega,\qquad \Omega=\{w\in K:F_i(w)\le b_i,\ i=1,2,3\}.
\end{equation}
Cách phát biểu này đòi hỏi cả ba tiêu chí đồng thời đạt mức chấp nhận được. Nó bỏ số hạng Tikhonov và không bảo toàn bài toán tối ưu gốc nào cũng như các nghiệm của chúng.

\begin{proposition}\label{prop:fs-reg}
Nếu $Q_R\succeq0$ và $\min_k\rho_k>0$ thì cả ba hàm trong~\eqref{eq:fsF1}--\eqref{eq:fsF3} trơn và lồi trên một lân cận mở lồi của $K$. Do đó~\eqref{eq:nsfp-fs} là một NSFP với hộp chuẩn mà mọi ràng buộc đều lồi. Vì vậy~\eqref{eq:nsfp-fs} thỏa (A1) và (A2) của Giả thiết~\ref{ass:standing}, với $U=U_\rho$, bất cứ khi nào $\Omega\ne\varnothing$.
\end{proposition}
\begin{proof}
Hai hàm đầu là bình phương chuẩn của các ánh xạ affine. Với $w\in U_\rho$ và một hướng $h$ bất kỳ, đặt $t=\rho^{\top}w$ và $a=\rho^{\top}h$. Lấy đạo hàm trực tiếp ta được
\begin{equation}\label{eq:fsF3-curvature}
D^2F_3(w)[h,h]=\frac2t\left(h-\frac atw\right)^{\!\top}Q_R\left(h-\frac atw\right)\ge0.
\end{equation}
Vậy $F_3$ lồi và trơn trên nửa không gian mở $U_\rho$. Vì $\rho^{\top}w\ge\min_k\rho_k>0$ với $w\in K$, nửa không gian này chứa $K$.
\end{proof}

Thực nghiệm dùng $I^{\cvx}=\{1,2\}$, $I^{\ps}=\{3\}$, một cách gán hợp lệ vì hàm lồi là giả lồi; khi đó DTAP chỉ dùng thông tin về dấu cho $F_3$, nên không có ràng buộc thực sự không lồi nào được kiểm tra. Cách gán toàn lồi $I^{\cvx}=\{1,2,3\}$ cũng hợp lệ, và Định lý~\ref{thm:strong} áp dụng cho cả hai. Mệnh đề~\ref{prop:fs-reg} cần $\min_k\rho_k>0$; trên mọi bài toán dùng ở đây, điểm số nhỏ nhất là $4.77\times10^{-5}$. Do đó ngưỡng bảo vệ $\max\{\rho^\top w,10^{-15}\}$ trong mã cài đặt không bao giờ có tác dụng, và $F_3$ được cài đặt đúng là~\eqref{eq:fsF3}.

Các ngưỡng được hiệu chỉnh trên dữ liệu huấn luyện: $b_i$ là phân vị thứ $q_i$ của $F_i$ trên $3\,000$ điểm Dirichlet (seed $42$) và $p$ đỉnh của $K$, trong đó một quy tắc cố định chọn $q_i$ theo cỡ mẫu, số đặc trưng, số lớp và mức mất cân bằng lớp (siêu dữ liệu bổ sung). Mọi bộ giải dùng chung các ngưỡng này. Hiệu chỉnh riêng từng $b_i$ không bảo đảm $\Omega\ne\varnothing$; Mục~\ref{ssec:feasibility-audit} kiểm tra tính chấp nhận được bằng số. Nếu ngưỡng được cho trước thì sẽ được dùng nguyên như vậy.

Với mọi phương pháp, sắp xếp các trọng số tính được thành $w^\star_{(1)}\ge\cdots\ge w^\star_{(p)}$ và giữ lại $s$ đặc trưng đầu, trong đó
\begin{equation}\label{eq:thetaRule}
s=\min\left\{s'\in\{1,\ldots,p\}:\sum_{j=1}^{s'}w^\star_{(j)}\ge\theta\right\},\qquad 0<\theta<1.
\end{equation}
Quy tắc trọng số tích lũy này không áp đặt một số lượng đặc trưng cho trước. Việc cắt bớt các trọng số không bảo toàn bảo đảm tính chấp nhận được của vector đầy đủ.

\begin{remark}\label{rem:fs-cost}
Với các ma trận đã tính sẵn, tính $F_1,\nabla F_1$ tốn $O(nrp)$, tính $F_2,\nabla F_2$ tốn $O(np)$, và tính $F_3,\nabla F_3$ tốn $O(p^2)$. Phép chiếu lên $K=\Delta_{p-1}$ bằng sắp xếp tốn $O(p\log p)$ phép toán~\citep{condat2016simplex}. Vì vậy một vòng lặp dày tốn $O(nrp+p^2)$ phép toán và cần $O(nrp+p^2)$ bộ nhớ, kể cả $Q_R$. Tiền xử lý và huấn luyện bộ phân loại là các chi phí riêng.
\end{remark}

\subsection{Thiết lập thực nghiệm}\label{ssec:setup}
Nghiên cứu phân loại chính dùng tám bộ dữ liệu UCI~\citep{kelly2023uci} trong Bảng~\ref{tab:datasets}. Các lần chạy DTAP báo cáo ở đây xuất phát từ vector trọng số đều và dùng $\rho_{\cvx}=1.98$, $\delta=10^{-3}$, $\gamma_k=(k+0.5)^{-0.55}$, sai số cho phép $10^{-6}$, và tối đa $1000$ vòng lặp. SG-const dùng bước $\nu_k=10^{-3}$ và SG-dyn dùng mô-đun lấy mẫu $\widehat L$ (Mục~\ref{sssec:baselines}); mô-đun tinh chỉnh chỉ được xét ở Mục~\ref{ssec:feasibility-audit}. N-Fscore dùng $\varepsilon_{\mathrm{PNN}}=10^{-3}$ và $1000$ điểm lưới Euler trên $[0,0.2]$. AFRSFCV dùng tốc độ học $10^{-3}$, $\lambda_1=1$, $\lambda_2=0.5$, và $1000$ bước. Ngưỡng tích lũy dùng chung là $\theta=0.9$.

\begin{table}[htbp]
\centering
\caption{Các bộ dữ liệu dùng cho phép so sánh phân loại chính. Bộ Movement Libras bổ sung chỉ được dùng trong phần chẩn đoán bộ giải.}\label{tab:datasets}
\begin{tabular}{lrrrl}
\toprule
Bộ dữ liệu   & Số mẫu $n$ & Số đặc trưng $p$ & Số lớp $r$ & Cân bằng    \\
\midrule
Glass        & $214$       & $9$          & $6$         & Mất cân bằng \\
Diabetes     & $768$       & $8$          & $2$         & Mất cân bằng \\
Segmentation & $2\,310$    & $18$         & $7$         & Cân bằng   \\
Abalone      & $4\,177$    & $8$          & $3$         & Cân bằng   \\
Satimage     & $6\,435$    & $36$         & $6$         & Mất cân bằng \\
Yeast        & $1\,484$    & $7$          & $10$        & Mất cân bằng \\
Spambase     & $4\,601$    & $57$         & $2$         & Mất cân bằng \\
Pendigits    & $10\,992$   & $16$         & $10$        & Cân bằng   \\
\bottomrule
\end{tabular}
\end{table}
\subsubsection{Các bộ giải cơ sở}\label{sssec:baselines}
Các đoạn sau mô tả các bộ giải cơ sở: N-Fscore và AFRSFCV cho các mô hình ở Mục~\ref{sssec:nfscore} và~\ref{sssec:afrsfcv-model}, và các phương pháp dưới gradient SG-const và SG-dyn.

\paragraph*{N-Fscore: mạng nơ-ron chiếu}
Để giải bài toán quy hoạch phân thức~\eqref{eq:wangOrig}, Wang và cộng sự~\citep{wang2021fractional} đề xuất một phương pháp động lực học nơ-ron ghép một mạng nơ-ron chiếu (PNN) với một mạng $k$-winners-take-all. PNN là hệ vi phân
\begin{equation}\label{eq:pnn}
\begin{split}\varepsilon_{\mathrm{PNN}}\frac{du}{dt}={}&-\Pi g(u)+\tfrac1p\mathbf{1}_p\\&-(I-\Pi)\bigl(u-g(u)+\nabla f((I-\Pi)g(u)+\tfrac1p\mathbf{1}_p)\bigr).\end{split}
\end{equation}
với đầu ra $w=g(u)$, trong đó $g(u)=\max\{0,u\}$ là hàm ReLU theo từng thành phần, $\Pi=\tfrac1p\mathbf{1}_p\mathbf{1}_p^{\top}$, $\nabla f$ là gradient của hàm mục tiêu trong~\eqref{eq:wangOrig}, và $\varepsilon_{\mathrm{PNN}}>0$ là một hằng số thời gian. Ta tích phân~\eqref{eq:pnn} bằng phương pháp Euler hiện trên một lưới đều gồm $N_{\mathrm{grid}}$ điểm trong $[0,T]$, không có quy tắc dừng, với các tham số ở Mục~\ref{ssec:setup}. Vector trọng số là hình chiếu $P_K\bigl(g(u(T))\bigr)$ của đầu ra cuối cùng của mạng lên $K$. Để so sánh đồng nhất, ta chỉ dùng phần PNN để thu được $w^\star$; giai đoạn $k$-WTA được thay bằng quy tắc ngưỡng tích lũy dùng chung ở Mục~\ref{ssec:selection}.

\paragraph*{AFRSFCV: gradient và bước cập nhật}
Hàm mất mát softmax~\eqref{eq:afrsfcvLoss} được cực tiểu hóa theo $\tau\in\R^p$ bằng hạ gradient. Gradient của nó là
\begin{equation}\label{eq:gradTau}
\nabla_\tau L(\tau)=\bigl(\diag(w(\tau))-w(\tau)w(\tau)^{\top}\bigr)\nabla_w L(w(\tau)),
\end{equation}
với $\nabla_w L=-\mathbf{M}^{\top}(\vecop(Y)-\mathbf{M}w)+\lambda_1(P^{\top}Pw-P^{\top}\mathbf{1}_n)+\lambda_2 w$, trong đó $\mathbf{M}\in\R^{nr\times p}$ có cột thứ $k$ là $\vecop(\underline{R_k}(D))$. Bước cập nhật $\tau^{t+1}=\tau^t-\eta\,\nabla_\tau L(\tau^t)$ chạy cố định $T$ bước; trọng số cuối $w^\star=w(\tau^T)\in K$ được đưa vào quy tắc chọn dùng chung ở Mục~\ref{ssec:selection}.

\paragraph*{SG-const và SG-dyn}
Các phương pháp dưới gradient cho bài toán chấp nhận tựa lồi này~\citep{hu2022quasi} theo
chiến lược ràng buộc bị vi phạm nhiều nhất của bài báo đó: với $\mathcal I_k=\argmax_i v_i(w^k)$,
\[
w^{k+1}=P_K(w^k-\nu_k\bar h_k),\qquad
\bar h_k=\frac1{|\mathcal I_k|}\sum_{i\in\mathcal I_k}\frac{\nabla F_i(w^k)}{\|\nabla F_i(w^k)\|}.
\]
SG-const dùng $\nu_k=10^{-3}$; SG-dyn dùng bước động $\nu_k=J(w^k)/L$ (số mũ H\"older
bằng $1$), với $L=\widehat L$ là chuẩn gradient lớn nhất tại $500$ điểm lấy mẫu trong $K$.



\subsubsection{Quy trình đánh giá}\label{sssec:protocol}
Mỗi bộ dữ liệu dùng một phép chia phân tầng duy nhất $80/20$ thành tập huấn luyện và tập kiểm tra
(seed $42$). Chuẩn hóa min--max được xác định một lần trên toàn bộ tập huấn luyện, không
xác định riêng trong từng fold chọn đặc trưng, rồi áp dụng cho dữ liệu kiểm tra.
Các fold phân tầng (seed $42$) trong tập huấn luyện dùng $N_f=10$, giảm còn $7$ trên Glass
và $4$ trên Yeast. Mặt nạ của mỗi fold dùng $\theta=0.9$ trong~\eqref{eq:thetaRule};
quy tắc đa số chặt cố định giữ lại các đặc trưng được chọn trong ít nhất
$\lfloor N_f/2\rfloor+1$ fold. Ta ký hiệu tập đặc trưng được bỏ phiếu chọn là $A^\star$. Các bộ phân loại cuối dùng toàn bộ tập huấn luyện với các
đặc trưng được chọn này: RBF-SVM ($C=1$, \texttt{gamma=scale}) và $5$-NN, với tập
đặc trưng đầy đủ làm đối chứng. Ta báo cáo độ chính xác và AUC macro theo cách một-chọi-phần-còn-lại (one-vs-rest).

Các mặt nạ chọn theo bỏ phiếu được tính một lần rồi dùng lại: giai đoạn đánh giá chỉ huấn luyện
các bộ phân loại phía sau trên các mặt nạ cố định này, với seed hiệu chỉnh xác suất của SVM
bằng $42$; các mặt nạ giống nhau cho cùng dự đoán. Dữ liệu kiểm tra không được dùng để chọn đặc trưng hay chọn bất kỳ tham số hoặc ngưỡng bỏ phiếu nào.
Phép chia huấn luyện/kiểm tra được cố định trước khi tính các tập đặc trưng, và quy tắc chọn
mức ngưỡng chỉ dùng cỡ mẫu, số đặc trưng, số lớp
và kích thước các lớp của toàn bộ dữ liệu, không dùng giá trị đặc trưng nào. Vì chỉ dùng một phép chia,
đây là một đánh giá khám phá, không phải một đánh giá đầu-cuối (end-to-end) hay kiểm định chéo lồng (nested).
Với ma trận chọn theo fold $Z\in\{0,1\}^{N_f\times p}$, chỉ số Nogueira~\citep{nogueira2018}
được tính bởi
\[
\Psi=1-\frac{p^{-1}\sum_{j=1}^{p}s_j^2}{\bar q(1-\bar q)},\quad
s_j^2=\frac{1}{N_f-1}\sum_{b=1}^{N_f}(Z_{bj}-\bar Z_j)^2,\quad
\bar q=\frac{1}{N_f p}\sum_{b,j}Z_{bj}.
\]
Ở đây $\bar Z_j=N_f^{-1}\sum_b Z_{bj}$, và $s_j^2$ là phương sai không chệch trên $N_f$
fold.
Vì các fold huấn luyện chồng lên nhau, độ ổn định chỉ được báo cáo ở mức mô tả.
Với mỗi độ đo phân loại, ta dùng kiểm định hạng có dấu hai phía trên
tám bộ dữ liệu, bỏ các hiệu bằng không, gán hạng trung bình cho các giá trị bằng nhau, và
liệt kê mọi cách đặt dấu của các hiệu khác không; các cặp toàn số 0 nhận giá trị $p$ bằng $1$.
Hiệu chỉnh Holm áp dụng riêng cho bốn phép so sánh DTAP với phương pháp cơ sở trong
mỗi độ đo. Kiểm định Friedman so sánh năm phương pháp chọn đặc trưng.

\subsection{Kết quả phân loại}\label{ssec:results}
Các Bảng~\ref{tab:nfeat}--\ref{tab:perf-summary} báo cáo kích thước tập được chọn,
độ ổn định và hiệu quả dự đoán. Trung bình DTAP giữ lại $84.2\%$ số đặc trưng,
trong đó $33/36$ trên Satimage và $49/57$ trên Spambase.
Độ ổn định trung bình của DTAP là $0.786$, so với $0.822$ của N-Fscore.
Như vậy DTAP giảm số đặc trưng ở mức vừa phải; nó không chọn được các tập thưa hơn hay
ổn định hơn các mô hình gốc. AFRSFCV chọn các tập nhỏ hơn nhiều trên một số bộ dữ liệu,
với độ chính xác trung bình thấp hơn trong cấu hình cố định.

Độ chính xác trung bình của DTAP cách đối chứng dùng toàn bộ đặc trưng không quá $0.44$ điểm phần trăm
với cả hai bộ phân loại. Nó cao hơn N-Fscore và AFRSFCV lần lượt $3.12$
và $4.83$ điểm về độ chính xác SVM, trong khi giữ lại nhiều đặc trưng hơn.
Đây là sự đánh đổi giữa kích thước và khả năng dự đoán, không phải một phép so sánh cùng số lượng đặc trưng.
Cả hai biến thể SG có mặt nạ chọn theo bỏ phiếu giống hệt DTAP trên bảy bộ dữ liệu;
chỉ Spambase khác ($48$ so với $49$ đặc trưng).
SG-dyn chạy hết 1000 vòng lặp tối đa ở 6/7 fold chọn trên Glass,
7/10 trên Diabetes và 10/10 trên Abalone, còn SG-const ở 1/10 trên Abalone
(DTAP không có fold nào); trọng số dùng để bỏ phiếu của cả hai biến thể SG là các điểm lặp tốt nhất
tìm được trong giới hạn số vòng lặp. Kết quả SG trên Spambase, bộ dữ liệu duy nhất
mà chúng khác DTAP, được ghi dưới Bảng~\ref{tab:acc};
Bảng~\ref{tab:perf-summary} gồm cả năm phương pháp.
Độ ổn định theo fold có thể khác nhau ngay cả khi mặt nạ chọn theo bỏ phiếu cuối cùng giống hệt.

Bảng~\ref{tab:wilcoxon} báo cáo các kiểm định hạng có dấu. Không phép so sánh từng cặp nào
bác bỏ ở mức $0.05$ sau hiệu chỉnh Holm; các phép so sánh SVM với
hai mô hình gốc có giá trị $p$ hiệu chỉnh trong khoảng $0.0625$ đến $0.0938$.
Các kiểm định Friedman tiệm cận bác bỏ giả thuyết bằng nhau trên tổng thể, nhưng điều này không xác lập
ưu thế của một cặp cụ thể nào. Mọi giá trị $p$ của các cặp DTAP--SG đều bằng $1$;
điều này phản ánh các giá trị bằng nhau và bằng chứng hạn chế, không phải một khẳng định tương đương thống kê.
Các phép so sánh trên mẫu nhỏ này liên quan đến tám bộ dữ liệu không đồng nhất và không
đo sự biến động khi lặp lại phép chia huấn luyện/kiểm tra.

\begin{table}[htbp]
\centering\small\setlength{\tabcolsep}{4pt}
\caption{Số đặc trưng được chọn $|A^\star|$ / độ ổn định Nogueira $\Psi$; $p$ là số chiều gốc. Cả năm phương pháp chọn đặc trưng đều được trình bày.}\label{tab:nfeat}
\begin{tabular}{lrrrrrr}
\toprule
Bộ dữ liệu & $p$ & N-Fscore & AFRSFCV & DTAP & SG-const & SG-dyn \\
\midrule
Glass & $9$ & $7/0.816$ & $7/0.609$ & $7/0.633$ & $7/0.633$ & $7/0.633$ \\
Diabetes & $8$ & $7/0.290$ & $3/0.810$ & $7/0.707$ & $7/0.707$ & $7/0.707$ \\
Segmentation & $18$ & $8/0.885$ & $2/1.000$ & $16/0.800$ & $16/0.888$ & $16/0.888$ \\
Abalone & $8$ & $4/0.950$ & $2/1.000$ & $5/1.000$ & $5/1.000$ & $5/0.870$ \\
Satimage & $36$ & $10/0.871$ & $3/0.685$ & $33/0.563$ & $33/0.563$ & $33/0.563$ \\
Yeast & $7$ & $6/1.000$ & $5/1.000$ & $6/0.757$ & $6/0.757$ & $6/0.757$ \\
Spambase & $57$ & $26/0.910$ & $4/0.885$ & $49/0.832$ & $48/0.851$ & $48/0.851$ \\
Pendigits & $16$ & $13/0.854$ & $12/0.063$ & $15/1.000$ & $15/1.000$ & $15/1.000$ \\
\bottomrule
\end{tabular}

\end{table}
\begin{table}[htbp]
\centering\small\setlength{\tabcolsep}{4pt}
\caption{Độ chính xác (\%) / AUC macro (one-vs-rest) theo từng bộ dữ liệu trên tập kiểm tra duy nhất. Các mặt nạ chọn theo bỏ phiếu giống nhau dùng chung dự đoán; các ngoại lệ của SG được ghi bên dưới.}\label{tab:acc}
\begin{tabular}{llrrrr}
\toprule
Bộ dữ liệu & Bộ phân loại & Đầy đủ & N-Fscore & AFRSFCV & DTAP \\
\midrule
Glass & SVM & $76.74/0.9222$ & $69.77/0.9073$ & $72.09/0.9174$ & $72.09/0.9174$ \\
 & $5$-NN & $72.09/0.9097$ & $74.42/0.8933$ & $69.77/0.9371$ & $69.77/0.9371$ \\
Diabetes & SVM & $74.68/0.8089$ & $71.43/0.7957$ & $65.58/0.6617$ & $74.68/0.8098$ \\
 & $5$-NN & $71.43/0.7356$ & $70.78/0.7154$ & $67.53/0.7145$ & $73.38/0.7766$ \\
Segmentation & SVM & $93.29/0.9936$ & $80.52/0.9714$ & $78.79/0.9668$ & $93.29/0.9940$ \\
 & $5$-NN & $94.37/0.9880$ & $87.01/0.9801$ & $86.58/0.9646$ & $94.59/0.9897$ \\
Abalone & SVM & $61.60/0.8102$ & $61.36/0.8065$ & $62.56/0.8114$ & $62.68/0.8123$ \\
 & $5$-NN & $59.81/0.7742$ & $61.00/0.7605$ & $58.73/0.7752$ & $60.29/0.7695$ \\
Satimage & SVM & $89.12/0.9796$ & $87.10/0.9756$ & $80.42/0.9611$ & $89.12/0.9797$ \\
 & $5$-NN & $91.14/0.9792$ & $87.80/0.9705$ & $80.96/0.9306$ & $91.22/0.9766$ \\
Yeast & SVM & $58.25/0.8561$ & $58.92/0.8241$ & $57.91/0.8207$ & $58.59/0.8426$ \\
 & $5$-NN & $51.52/0.7667$ & $51.52/0.7763$ & $51.52/0.7681$ & $51.52/0.7585$ \\
Spambase & SVM & $92.07/0.9662$ & $88.38/0.9392$ & $87.62/0.9155$ & $91.97/0.9616$ \\
 & $5$-NN & $90.34/0.9443$ & $85.78/0.9214$ & $88.49/0.9206$ & $89.25/0.9437$ \\
Pendigits & SVM & $99.09/0.9995$ & $98.95/0.9996$ & $97.77/0.9996$ & $98.95/0.9995$ \\
 & $5$-NN & $98.95/0.9982$ & $98.64/0.9982$ & $97.73/0.9965$ & $99.05/0.9987$ \\
\bottomrule
\end{tabular}
\par\smallskip\begin{minipage}{0.97\textwidth}\footnotesize Trên bảy bộ dữ liệu còn lại, cả hai biến thể SG cho kết quả đúng bằng DTAP. Trên Spambase, cả hai cho $91.86/0.9628$ (SVM) và $89.14/0.9441$ ($5$-NN). \end{minipage}
\end{table}
\begin{table}[htbp]
\centering\small\setlength{\tabcolsep}{4pt}
\caption{Hiệu quả phân loại trung bình không trọng số trên tám bộ dữ liệu phân loại. Gồm cả năm phương pháp chọn đặc trưng và đối chứng dùng toàn bộ đặc trưng.}\label{tab:perf-summary}
\begin{tabular}{lrrrr}
\toprule
Phương pháp & Độ chính xác SVM & Độ chính xác $5$-NN & AUC SVM & AUC $5$-NN \\
\midrule
Đầy đủ & $80.61$ & $78.71$ & $0.9170$ & $0.8870$ \\
N-Fscore & $77.05$ & $77.12$ & $0.9024$ & $0.8770$ \\
AFRSFCV & $75.34$ & $75.16$ & $0.8818$ & $0.8759$ \\
DTAP & $80.17$ & $78.63$ & $0.9146$ & $0.8938$ \\
SG-const & $80.16$ & $78.62$ & $0.9148$ & $0.8939$ \\
SG-dyn & $80.16$ & $78.62$ & $0.9148$ & $0.8939$ \\
\bottomrule
\end{tabular}

\end{table}
\begin{table}[htbp]
\centering\footnotesize\setlength{\tabcolsep}{4pt}
\caption{DTAP so với từng phương pháp cơ sở: giá trị $p$ của kiểm định hạng có dấu, gốc / sau hiệu chỉnh Holm. Mỗi độ đo có bốn phép so sánh trên tám bộ dữ liệu. Cột cuối là giá trị $p$ tiệm cận của kiểm định Friedman (đã hiệu chỉnh cho giá trị bằng nhau) cho cả năm phương pháp.}\label{tab:wilcoxon}
\begin{tabular}{lrrrrr}
\toprule
Độ đo & N-Fscore & AFRSFCV & SG-const & SG-dyn & Friedman \\
\midrule
Độ chính xác, SVM & $0.0312/0.0938$ & $0.0156/0.0625$ & $1.0000/1.0000$ & $1.0000/1.0000$ & $0.0005$ \\
Độ chính xác, $5$-NN & $0.3750/1.0000$ & $0.0312/0.1250$ & $1.0000/1.0000$ & $1.0000/1.0000$ & $0.0057$ \\
AUC, SVM & $0.0156/0.0625$ & $0.0312/0.0938$ & $1.0000/1.0000$ & $1.0000/1.0000$ & $0.0055$ \\
AUC, $5$-NN & $0.0781/0.3125$ & $0.1562/0.4688$ & $1.0000/1.0000$ & $1.0000/1.0000$ & $0.0435$ \\
\bottomrule
\end{tabular}

\end{table}

\subsection{Mức hoạt động của ràng buộc và một phép so sánh với phương pháp dưới gradient đã tinh chỉnh}
\label{ssec:feasibility-audit}

Như đã nêu ở Mục~\ref{ssec:reduction}, các ngưỡng đã hiệu chỉnh không nhất thiết cho
$\Omega\ne\varnothing$. Với mỗi bài toán dữ liệu đầy đủ dưới đây, ta giữ lại điểm đã tính có độ chênh
$\min_i(b_i-F_i(w))$ lớn nhất trên mọi lần chạy. Khi tính lại trên các ma trận dựng lại từ dữ liệu, mỗi
điểm này thỏa chặt cả ba bất đẳng thức; độ chênh nhỏ nhất là
$4.33\times10^{-4}$. Đây là một phép kiểm tra số, không phải chứng nhận chính xác về tính chấp nhận được.

Ta gọi các điểm lặp của một lần chạy, kể cả điểm cuối, là các điểm đã tính của lần chạy đó, và đếm các điểm tại đó
$F_3(w^k)-b_3>10^{-6}$. Đại lượng này đo sự vi phạm dọc theo quỹ đạo, chứ không đo việc
ràng buộc có chặt tại điểm giới hạn hay không. Bảng~\ref{tab:activation-tuning} cho thấy, từ
điểm xuất phát chuẩn, $F_3$ không bị vi phạm tại điểm đã tính nào trên bảy trong tám bộ dữ liệu
phân loại, và bị vi phạm tại một trong $48$ điểm đã tính trên Abalone. Vì $F_3$ lồi
(Mệnh đề~\ref{prop:fs-reg}), các bài toán chọn đặc trưng này kiểm tra một lớp con lồi và
không cho thấy cần đến phần mở rộng không lồi. Ta cũng báo cáo phương án so sánh toàn lồi
(AC), gán cả ba ràng buộc vào nhánh lồi.

\paragraph*{Quy trình và chi phí tinh chỉnh.}
Đây là các phép thử thuần túy về tối ưu trên các bài toán chấp nhận với dữ liệu đầy đủ.
Chúng dùng $J\le10^{-6}$, tối đa $1\,000$ điểm đã tính, và cùng
điểm xuất phát cho mọi phương pháp. SG-const và SG-dyn được định nghĩa ở Mục~\ref{sssec:baselines}. Bên cạnh SG-dyn với
mô-đun lấy mẫu $\widehat L$, một biến thể tinh chỉnh chọn $L=c\widehat L$ từ lưới đặt thủ công
$c\in\{10^{-3},10^{-2.5},\ldots,10^1\}$ bằng ba điểm xuất phát
ngẫu nhiên (seed $1101$--$1103$). Việc chọn trước hết cực đại hóa số
lần chạy thành công, sau đó cực tiểu hóa tổng số lần tính hàm, rồi tổng phần dư tốt nhất.
Đánh giá dùng điểm xuất phát chuẩn và ba điểm xuất phát ngẫu nhiên khác
(seed $2101$--$2103$). Các điểm xuất phát cho bài toán đặc trưng là các mẫu Dirichlet.
Chỉ các điểm xuất phát được tách riêng: tinh chỉnh và đánh giá dùng cùng
bài toán dữ liệu đầy đủ, và không dùng kết quả phân loại nào. Cả
$\widehat L$ lẫn giá trị đã tinh chỉnh đều không phải là một mô-đun H\"older toàn cục đã được chứng nhận.

\begin{table}[t]
\centering\small
\setlength{\tabcolsep}{4pt}
\caption{Số lần tính hàm trong tối ưu, xuất phát từ điểm đều trên đơn hình.
AC là phương án so sánh toàn lồi; ``lấy mẫu'' và ``tinh chỉnh'' là hai biến thể của SG-dyn.
$\dagger$ nghĩa là không đạt sai số cho phép trong giới hạn số điểm.
Cột cuối là số điểm mà DTAP vi phạm ràng buộc thứ ba, chia cho
số điểm đã tính của nó.}
\label{tab:activation-tuning}
\begin{tabular}{@{}lrrrrrl@{}}
\toprule
Bộ dữ liệu & DTAP & AC & SG-const & Lấy mẫu & Tinh chỉnh & Hoạt động của $F_3$\\
\midrule
Glass        & 89 & 89 & 520 & $1000^\dagger$ & 100 & $0/89$\\
Diabetes     & 14 & 14 & 350 & $1000^\dagger$ & 51 & $0/14$\\
Segmentation & 3 & 3 & 124 & 170 & 3 & $0/3$\\
Abalone      & 48 & 48 & 638 & $1000^\dagger$ & 3 & $1/48$\\
Satimage     & 2 & 2 & 60 & 102 & 2 & $0/2$\\
Yeast        & 4 & 4 & 324 & 305 & 3 & $0/4$\\
Spambase     & 5 & 5 & 93 & 295 & 56 & $0/5$\\
Pendigits    & 2 & 2 & 49 & 10 & 31 & $0/2$\\
\midrule
Movement Libras & 637 & 637 & 62 & $1000^\dagger$ & 144 & $0/637$\\
\bottomrule
\end{tabular}
\end{table}

Mỗi khi $F_3(w^k)\le b_3$ dọc theo một lần chạy, tập active của nhánh giả lồi rỗng,
$\mu_k=1$ và DTAP trùng từng bước với phương án so sánh toàn lồi;
trên tám bài toán có mức hoạt động của $F_3$ bằng không, hai cột DTAP và AC giống hệt nhau
phản ánh sự trùng khớp này. Trên Abalone, $F_3$ chỉ bị vi phạm tại điểm xuất phát
đều ($1/48$), nơi $J^{\ps}_0\approx1.56\times10^{-3}$ rất nhỏ so với
$J^{\cvx}_0\approx366.3$, nên $1-\mu_0\approx4.27\times10^{-6}$; các giá trị của $J$
dọc hai quỹ đạo thỏa $|J_{\rm DTAP}-J_{\rm AC}|\le3.18\times10^{-5}\,J_{\rm AC}$ tại
mọi điểm đã tính,
và cả hai lần chạy dừng sau $48$ điểm đã tính. Vì vậy các cột này
không phải là phép so sánh hai cách gán; hai cách gán chỉ khác nhau về số lần thành công
trên các điểm xuất phát ngẫu nhiên ($26/36$ so với $36/36$).

Việc tinh chỉnh nâng số lần thành công của SG-dyn, xuất phát từ điểm chuẩn, từ $5/9$ lên $9/9$.
SG-dyn tinh chỉnh có thể tốt hơn DTAP, chẳng hạn trên Abalone, và SG-const nhanh hơn
trên Movement Libras. Trên cả bốn điểm xuất phát của cả chín bài toán đặc trưng,
số lần thành công là $26/36$ cho DTAP, $36/36$ cho phương án so sánh toàn lồi,
$33/36$ cho SG-const, $15/36$ cho SG-dyn lấy mẫu, và $28/36$
cho SG-dyn tinh chỉnh. Một biến thể DTAP cân bằng cũng đạt $26/36$; nó gán
trọng số bằng nhau khi cả hai nhánh bị vi phạm, và ngược lại gán toàn bộ trọng số cho nhánh duy nhất bị vi phạm.
Như vậy các phép thử này không chứng tỏ được ưu thế đồng đều của trọng số tỉ lệ,
cũng không chứng tỏ được tính ổn định đối với điểm khởi đầu.

Chín lần tìm kiếm tham số trên các bài toán đặc trưng cần tổng cộng $179\,744$ lần tính hàm trong tối ưu
và $10.80$ giây thời gian giải cho việc tinh chỉnh. Các chi phí này
là phần thêm vào so với Bảng~\ref{tab:activation-tuning}; DTAP không được tinh chỉnh.
Việc lấy mẫu mô-đun ban đầu, tiền xử lý và chi phí phụ của việc tìm kiếm không được tính
vào tổng thời gian giải này. Hồ sơ bổ sung ghi chi phí của
từng lần tìm kiếm.

\paragraph*{Một ràng buộc giả lồi thực sự không lồi.}
Để kiểm tra trực tiếp nhánh giả lồi, xét $K=[-2,2]^d$ và
\begin{equation}\label{eq:logball-test}
\begin{aligned}
 F_1(x)&=\|x+e_1\|^2\le1,\qquad
 F_2(x)=\|x-e_2\|^2\le\tfrac94,\\
 F_3(x)&=\kappa\log(1+\|x-e_1\|^2)\le \kappa\log3.
\end{aligned}
\end{equation}
Ở đây $d\in\{2,10,50\}$ và $\kappa\in\{0.3,1,3\}$ cho chín cấu hình.
Hàm thứ ba là giả lồi: nó là hợp của một hàm bậc hai lồi với một hàm khả vi tăng
có đạo hàm dương. Nó không
lồi, vì giá trị riêng theo phương bán kính của ma trận Hesse là
$2\kappa(1-r^2)/(1+r^2)^2<0$ với $r=\|x-e_1\|>1$.
Điểm minh chứng chính xác $x=-e_1/10$ có các độ chênh dương
$19/100$, $31/25$ và $\kappa\log(300/221)$.
Điểm xuất phát chuẩn là $1.5\mathbf1$; ba điểm xuất phát dùng để đánh giá và ba
điểm xuất phát dùng để tinh chỉnh được lấy đều trên $[0.8,1.8]^d$ với các seed nêu trên
(NumPy \texttt{default\_rng}).
Để lấy mẫu mô-đun, ta dùng $500$ điểm đều trong $K$ (seed $42$).
SG-const dùng cùng bước $\nu_k=10^{-3}$ như trên các bài toán đặc trưng. SG-const
tinh chỉnh dùng một bước hằng $\nu$ chọn từ lưới
$\{10^{-3},10^{-2.5},\ldots,10^1\}$ với các điểm xuất phát tinh chỉnh và quy tắc chọn của
SG-dyn tinh chỉnh; nó chọn $\nu=10^{-0.5}$ ở cả chín cấu hình. Biến thể cân bằng
là biến thể đã mô tả ở trên.

\begin{table}[t]
\centering\small
\setlength{\tabcolsep}{5pt}
\caption{Các phép thử giả lồi không lồi, nhóm theo $\kappa$.
Mỗi ô là số lần thành công trên $12$ lần chạy, kèm trung vị số
điểm đã tính trong ngoặc, kể cả các lần chạy hết số điểm tối đa. Mỗi nhóm
gồm cả ba số chiều và bốn điểm xuất phát. SG-const tinh chỉnh dùng bước
được chọn như mô tả trong bài.}
\label{tab:synthetic-scale}
\begin{tabular}{@{}lrrrrrr@{}}
\toprule
$\kappa$ & DTAP & Cân bằng & SG-const & SG-const tinh chỉnh & SG-dyn lấy mẫu & SG-dyn tinh chỉnh\\
\midrule
$0.3$ & $0\ (1000)$ & $0\ (1000)$ & $0\ (1000)$ & $12\ (15)$ & $8\ (725.5)$ & $12\ (142.5)$\\
$1$   & $12\ (70)$ & $12\ (395.5)$ & $0\ (1000)$ & $12\ (15)$ & $12\ (301.5)$ & $12\ (44)$\\
$3$   & $12\ (4)$ & $12\ (4)$ & $0\ (1000)$ & $12\ (16)$ & $12\ (167)$ & $12\ (11)$\\
\bottomrule
\end{tabular}
\end{table}

DTAP gặp vi phạm ràng buộc thứ ba ở mọi lần trong $36$
lần chạy này. Tính gộp lại, ràng buộc thứ ba bị vi phạm tại $12\,835/12\,877$ điểm đã tính.
Tuy vậy, nó chỉ đạt sai số cho phép ở $24/36$ lần chạy, so
với $32/36$ của SG-dyn lấy mẫu, $36/36$ của SG-dyn tinh chỉnh và $36/36$ của SG-const
tinh chỉnh (Bảng~\ref{tab:synthetic-scale}).
SG-const tinh chỉnh đạt sai số cho phép ở mọi lần chạy với trung vị $15$--$16$
điểm đã tính, ít hơn mọi phương pháp khác tại $\kappa=0.3$ và $\kappa=1$; tại
$\kappa=3$, DTAP và biến thể cân bằng cần trung vị $4$. Các số đếm này
không gồm $147\,331$ lần tính hàm trong việc tìm kiếm tham số của nó.

Với $\nu_k=10^{-3}$, SG-const di chuyển nhiều nhất $999\times10^{-3}=0.999$ so với điểm xuất phát
trong giới hạn số điểm, trong khi mọi điểm xuất phát cách
$\Omega$ ít nhất $1.11$; vì vậy các số đếm bằng không của nó do cỡ bước và giới hạn số điểm quyết định.
Với $\kappa\le1$, nhánh lồi được thỏa từ điểm đã tính thứ hai trở đi,
nên sau đó DTAP đi theo phép lặp một nhánh của Hệ quả~\ref{cor:ps}; cả hai nhánh
cùng bị vi phạm chỉ tại $32$ trong $12\,877$ điểm đã tính này.

Mọi thất bại của DTAP đều xảy ra tại $\kappa=0.3$. Thay đổi $\kappa$ không làm thay đổi tập chấp nhận được,
nhưng $v_3$ và $\nabla F_3$ đều tỉ lệ với $\kappa$, nên khi
$\|g^{\ps}_k\|\le1$ bước giả lồi $(1-\mu_k)\gamma_k g^{\ps}_k$ tỉ lệ với
$\gamma_k\kappa^2$. Bảng~\ref{tab:dtap-sweep} lần lượt thay đổi từng tham số của DTAP
với $\kappa\in\{0.3,1\}$. Với $\gamma_k=a(k+0.5)^{-0.55}$, $a=10$ và $a=100$ cho
$12/12$ lần thành công tại $\kappa=0.3$; $\rho_{\cvx}=1$ cho $4/12$, đều với $d=2$;
$\delta=10^{-1}$ tái tạo chính xác các điểm lặp của cấu hình gốc (Nhận xét~\ref{rem:delta-role}); các
bước giảm nhanh hơn $s=0.75$ và $s=1$ không thành công lần nào tại $\kappa=0.3$, và
$s=1$ cũng thất bại cả $12$ lần chạy tại $\kappa=1$. Vì vậy các thất bại tại $\kappa=0.3$
phù hợp với việc bước giảm dần quá nhỏ so với thang của ràng buộc
này, chứ không phải với một thay đổi về tính chấp nhận được. Các phép thử này thay đổi
từng tham số một trên một họ bài toán và không chỉ ra được một cấu hình
tốt nhất nói chung.
Trên~\eqref{eq:logball-test}, $\delta=10^{-1}$ và $\delta=10^{-3}$ cho các dãy lặp giống hệt nhau vì
nhánh giả lồi chỉ có một chỉ số, tại mọi điểm xuất phát hoặc chỉ $F_1$ bị vi phạm hoặc
$v_1-v_2\ge4.79$, và nhánh lồi được thỏa từ điểm đã tính thứ hai trở đi; vì vậy
dòng $\delta$ của Bảng~\ref{tab:dtap-sweep} là một phép kiểm tra tính nhất quán, không phải bằng chứng về
vai trò của $\delta$.
Các lần tìm kiếm tham số trên bài toán tổng hợp tốn $188\,295$ lần tính hàm trong tối ưu và
$10.98$ giây thời gian giải cho SG-dyn, và $147\,331$ lần tính hàm và $8.27$
giây cho SG-const, với cùng các khoản loại trừ như trên, trong môi trường
nêu dưới đây.

\begin{table}[t]
\centering\small
\setlength{\tabcolsep}{5pt}
\caption{DTAP trên~\eqref{eq:logball-test} khi thay đổi từng tham số một,
$\gamma_k=a(k+0.5)^{-s}$. Cấu hình gốc $\rho_{\cvx}=1.98$, $\delta=10^{-3}$, $a=1$,
$s=0.55$ là cột DTAP của Bảng~\ref{tab:synthetic-scale}. Mỗi ô là
số lần thành công trên $12$ lần chạy (ba số chiều, bốn điểm xuất phát), kèm
trung vị số điểm đã tính trong ngoặc, kể cả các lần chạy hết số điểm tối đa.}
\label{tab:dtap-sweep}
\begin{tabular}{@{}lrr@{}}
\toprule
Cấu hình & $\kappa=0.3$ & $\kappa=1$\\
\midrule
Cấu hình gốc & $0\ (1000)$ & $12\ (70)$\\
$\rho_{\cvx}=1$ & $4\ (1000)$ & $12\ (223.5)$\\
$\delta=10^{-1}$ & $0\ (1000)$ & $12\ (70)$\\
$a=10$ & $12\ (77)$ & $12\ (6)$\\
$a=100$ & $12\ (6)$ & $12\ (48)$\\
$s=0.75$ & $0\ (1000)$ & $12\ (374)$\\
$s=1$ & $0\ (1000)$ & $0\ (1000)$\\
\bottomrule
\end{tabular}
\end{table}

\paragraph*{Khả năng tái lập.}
Kết quả chọn đặc trưng được lưu với Python~3.11.7 và NumPy~1.26.4
trên Windows, với yêu cầu một luồng BLAS/OpenMP. Các hàm tính giá trị (oracle) cho bài toán đặc trưng dùng dạng toàn phương
nén theo đại số; sai khác đã kiểm tra so với giá trị và gradient
gốc lần lượt không quá $1.82\times10^{-12}$ và
$1.10\times10^{-11}$. Thời gian giải báo cáo cho các bài toán đặc trưng ứng
với cách cài đặt nén này. Gói mã đi kèm chỉ chứa các hàm tính giá trị không nén;
nó tái tạo các số lần tính hàm, còn thời gian giải được báo cáo là của
cách cài đặt nén ban đầu. Các Bảng~\ref{tab:synthetic-scale}
và~\ref{tab:dtap-sweep} và chi phí tinh chỉnh trên bài toán tổng hợp được tạo với
Python~3.13.4 và NumPy~2.1.3 trên Windows với số luồng BLAS
mặc định; các số đếm của chúng là tất định, còn thời gian giải là các
phép đo một lần trên máy đó. Với Python~3.13.4 và NumPy~2.1.3, gói mã còn
tái tạo chính xác mọi giá trị số trong các Bảng~\ref{tab:datasets}--\ref{tab:wilcoxon}
(Bảng~\ref{tab:perf-summary} và~\ref{tab:wilcoxon} từ điểm chưa làm tròn) và trong
Bảng~\ref{tab:activation-tuning}, cùng các số đếm nêu ở Mục~\ref{ssec:results}
và~\ref{ssec:feasibility-audit}.
Các script bổ sung, quỹ đạo thô, hồ sơ tinh chỉnh, điểm minh chứng,
mã băm của mã nguồn và hướng dẫn kiểm tra được cung cấp kèm theo kết quả.

\section{Kết luận}\label{sec:conclusion}
DTAP giải các bài toán chấp nhận tách phi tuyến với hộp chuẩn và các thành phần lồi và giả lồi, bằng cách kết hợp một bước lồi điều chỉnh theo phần dư với một bước giả lồi có bảo vệ. Dưới giả thiết nhất quán và khả vi liên tục, toàn bộ dãy hội tụ tới một điểm chấp nhận được; không có mô-đun chính quy nào là đầu vào, mặc dù phân tích dùng các cận trên một tập compact chứa dãy lặp. Trong ứng dụng chọn đặc trưng, mọi ràng buộc đều lồi, và các bài toán thử có ràng buộc giả lồi không lồi kiểm tra nhánh giả lồi. Trong thực nghiệm, các phương pháp dưới gradient đã tinh chỉnh có thể bằng hoặc tốt hơn DTAP, và hiệu năng của DTAP với số vòng lặp hữu hạn phụ thuộc vào điểm xuất phát và thang cỡ bước. Tốc độ cho toàn bộ phần dư $J(x^k)$, ngoài cận của nhánh lồi ở Nhận xét~\ref{rem:rate}, và các bảo đảm dưới điều kiện cận sai số vẫn còn là vấn đề mở.

\section*{Tuyên bố về xung đột lợi ích}
Các tác giả tuyên bố không có xung đột lợi ích tài chính hay quan hệ cá nhân nào đã biết có thể ảnh hưởng đến công trình được báo cáo trong bài.
\section*{Tài trợ}
Các tác giả không nhận tài trợ riêng nào cho công trình này.
\section*{Tuyên bố về việc dùng AI tạo sinh}
OpenAI Codex đã được dùng trong quá trình chỉnh sửa để biên tập ngôn ngữ.
\section*{Tính sẵn có của dữ liệu}
Các bộ dữ liệu có tại UCI Machine Learning Repository~\citep{kelly2023uci}. Gói tái lập đi kèm chứa các script, kết quả số đã ghi lại, và hướng dẫn tạo lại các bảng đã báo cáo.

\section*{ORCID}
\noindent Nguyen Van The~\raisebox{-0.1ex}{\textit{https://orcid.org/0009-0004-6548-420X}}\\
Tran Ngoc Thang~\raisebox{-0.1ex}{\textit{https://orcid.org/0009-0000-9549-5736}}\\
Nguyen Hong Son~\raisebox{-0.1ex}{\textit{https://orcid.org/0009-0006-3106-8296}}


\begin{thebibliography}{99}
\newcommand{\bibdoi}[1]{doi: {\urlstyle{rm}\url{#1}}}

\bibitem[Bauschke and Combettes(2011)]{bauschke2011}
H.~H. Bauschke and P.~L. Combettes.
\newblock \emph{Convex Analysis and Monotone Operator Theory in {H}ilbert Spaces}.
\newblock CMS Books in Mathematics. Springer, New York, 2011.
\newblock \bibdoi{10.1007/978-1-4419-9467-7}.

\bibitem[Brown et~al.(2012)Brown, Pocock, Zhao, and Luj{\'a}n]{brown2012}
G.~Brown, A.~Pocock, M.-J.~Zhao, and M.~Luj{\'a}n.
\newblock Conditional likelihood maximization: a unifying framework for information theoretic feature selection.
\newblock \emph{Journal of Machine Learning Research}, 13:\penalty0 27--66, 2012.

\bibitem[Byrne(2002)]{byrne2002}
C.~Byrne.
\newblock Iterative oblique projection onto convex sets and the split feasibility problem.
\newblock \emph{Inverse Problems}, 18\penalty0 (2):\penalty0 441--453, 2002.
\newblock \bibdoi{10.1088/0266-5611/18/2/310}.

\bibitem[Censor and Elfving(1994)]{censor1994}
Y.~Censor and T.~Elfving.
\newblock A multiprojection algorithm using {B}regman projections in a product space.
\newblock \emph{Numerical Algorithms}, 8:\penalty0 221--239, 1994.
\newblock \bibdoi{10.1007/BF02142692}.

\bibitem[Combettes(2001)]{combettes2001quasifejer}
P.~L. Combettes.
\newblock Quasi-Fej{\'e}rian analysis of some optimization algorithms.
\newblock In D.~Butnariu, Y.~Censor, and S.~Reich, editors, \emph{Inherently Parallel Algorithms in Feasibility and Optimization and Their Applications}, volume~8 of \emph{Studies in Computational Mathematics}, pages 115--152. Elsevier, Amsterdam, 2001.
\newblock \bibdoi{10.1016/S1570-579X(01)80010-0}.

\bibitem[Condat(2016)]{condat2016simplex}
L.~Condat.
\newblock Fast projection onto the simplex and the $\ell_1$ ball.
\newblock \emph{Mathematical Programming}, 158\penalty0 (1):\penalty0 575--585, 2016.
\newblock \bibdoi{10.1007/s10107-015-0946-6}.

\bibitem[Hu et~al.(2022)Hu, Li, Yu, and Yip]{hu2022quasi}
Y.~Hu, G.~Li, C.~K.~W. Yu, and T.~L. Yip.
\newblock Quasi-convex feasibility problems: Subgradient methods and convergence rates.
\newblock \emph{European Journal of Operational Research}, 298\penalty0 (1):\penalty0 45--58, 2022.
\newblock \bibdoi{10.1016/j.ejor.2021.09.029}.

\bibitem[Kelly et~al.(2023)Kelly, Longjohn, and Nottingham]{kelly2023uci}
M.~Kelly, R.~Longjohn, and K.~Nottingham.
\newblock The {UCI} machine learning repository.
\newblock \url{https://archive.ics.uci.edu}, 2023.
\newblock Accessed 2026.

\bibitem[Li et~al.(2013)Li, Han, and Zhang]{li2013nsfp}
Z.~Li, D.~Han, and W.~Zhang.
\newblock A self-adaptive projection-type method for nonlinear multiple-sets split feasibility problem.
\newblock \emph{Inverse Problems in Science and Engineering}, 21\penalty0 (1):\penalty0 155--170, 2013.
\newblock \bibdoi{10.1080/17415977.2012.677445}.

\bibitem[L{\'o}pez et~al.(2012)L{\'o}pez, Mart{\'\i}n-M{\'a}rquez, Wang, and Xu]{lopez2012}
G.~L{\'o}pez, V.~Mart{\'\i}n-M{\'a}rquez, F.~Wang, and H.-K. Xu.
\newblock Solving the split feasibility problem without prior knowledge of matrix norms.
\newblock \emph{Inverse Problems}, 28\penalty0 (8):\penalty0 085004, 2012.
\newblock \bibdoi{10.1088/0266-5611/28/8/085004}.

\bibitem[Mangasarian(1965)]{mangasarian1965}
O.~L. Mangasarian.
\newblock Pseudo-convex functions.
\newblock \emph{SIAM Journal on Control}, 3\penalty0 (2):\penalty0 281--290, 1965.
\newblock \bibdoi{10.1137/0303020}.

\bibitem[Mangasarian(1994)]{mangasarian1994}
O.~L. Mangasarian.
\newblock \emph{Nonlinear Programming}.
\newblock Classics in Applied Mathematics. SIAM, Philadelphia, 1994.
\newblock \bibdoi{10.1137/1.9781611971255}.

\bibitem[Nogueira et~al.(2018)Nogueira, Sechidis, and Brown]{nogueira2018}
S.~Nogueira, K.~Sechidis, and G.~Brown.
\newblock On the stability of feature selection algorithms.
\newblock \emph{Journal of Machine Learning Research}, 18\penalty0 (174):\penalty0 1--54, 2018.

\bibitem[Polyak(1969)]{polyak1969}
B.~T. Polyak.
\newblock Minimization of unsmooth functionals.
\newblock \emph{USSR Computational Mathematics and Mathematical Physics}, 9\penalty0 (3):\penalty0 14--29, 1969.
\newblock \bibdoi{10.1016/0041-5553(69)90061-5}.

\bibitem[Wang et~al.(2021)Wang, Li, and Wang]{wang2021fractional}
Y.~Wang, X.~Li, and J.~Wang.
\newblock A neurodynamic optimization approach to supervised feature selection via fractional programming.
\newblock \emph{Neural Networks}, 136:\penalty0 194--206, 2021.
\newblock \bibdoi{10.1016/j.neunet.2021.01.004}.

\bibitem[Wang et~al.(2026)Wang, Zhang, An, and Deng]{wang2026afrsfcv}
C.~Wang, Y.~Zhang, S.~An, and T.~Deng.
\newblock Adaptive feature selection based on fuzzy rough set fusion model with class variance.
\newblock \emph{Pattern Recognition}, 170:\penalty0 112014, 2026.
\newblock \bibdoi{10.1016/j.patcog.2025.112014}.

\bibitem[Yang(2004)]{yang2004}
Q.~Yang.
\newblock The relaxed {CQ} algorithm solving the split feasibility problem.
\newblock \emph{Inverse Problems}, 20\penalty0 (4):\penalty0 1261--1266, 2004.
\newblock \bibdoi{10.1088/0266-5611/20/4/014}.

\bibitem[Yang et~al.(2014)Yang, Yang, and Zhang]{yang2014alternating}
Y.~Yang, Q.~Yang, and S.~Zhang.
\newblock Modified alternating direction methods for the modified multiple-sets split feasibility problems.
\newblock \emph{Journal of Optimization Theory and Applications}, 163\penalty0 (1):\penalty0 130--147, 2014.
\newblock \bibdoi{10.1007/s10957-013-0502-6}.

\bibitem[Yen et~al.(2019)Yen, Huyen, and Muu]{yen2019}
L.~H. Yen, N.~T.~T. Huyen, and L.~D. Muu.
\newblock A subgradient algorithm for a class of nonlinear split feasibility problems: application to jointly constrained Nash equilibrium models.
\newblock \emph{Journal of Global Optimization}, 73:\penalty0 849--868, 2019.
\newblock \bibdoi{10.1007/s10898-018-00735-0}.

\bibitem[Zhang et~al.(2009)Zhang, Han, and Li]{zhang2009adaptive}
W.~Zhang, D.~Han, and Z.~Li.
\newblock A self-adaptive projection method for solving the multiple-sets split feasibility problem.
\newblock \emph{Inverse Problems}, 25\penalty0 (11):\penalty0 115001, 2009.
\newblock \bibdoi{10.1088/0266-5611/25/11/115001}.

\bibitem[Zhang et~al.(2012)Zhang, Han, and Yuan]{zhang2012simultaneous}
W.~Zhang, D.~Han, and X.~Yuan.
\newblock An efficient simultaneous method for the constrained multiple-sets split feasibility problem.
\newblock \emph{Computational Optimization and Applications}, 52\penalty0 (3):\penalty0 825--843, 2012.
\newblock \bibdoi{10.1007/s10589-011-9429-8}.

\bibitem[Zhao and Yang(2011)]{zhao2011adaptive}
J.~Zhao and Q.~Yang.
\newblock Self-adaptive projection methods for the multiple-sets split feasibility problem.
\newblock \emph{Inverse Problems}, 27\penalty0 (3):\penalty0 035009, 2011.
\newblock \bibdoi{10.1088/0266-5611/27/3/035009}.

\end{thebibliography}

\end{document}
